% 第八章  补遗
\chapter{补遗}\label{ch:VIII}\index{buyi@补遗(complements)}

在本章中, 我们对类域论作一些补充, 并给出一些应用. 特别地, 我们把 \cite{Serre1970} 的第 I--IV、VI 章中证明的结果推广到任意数域上.

\section{局部的 \texorpdfstring{$n$}{n} 次幂何时是整体的 \texorpdfstring{$n$}{n} 次幂?}\label{sec:VIII.1}

显然, 若数域 $K$ 中的元素 $a$ 是 $K$ 中的 $n$ 次幂, 则它对所有素 $v$ 都是 $K_v$ 中的 $n$ 次幂. 本节研究其逆命题在何种意义上成立.

\begin{theorem}\label{thm:viii.1.1}
设 $K$ 是包含本原 $n$ 次单位根的数域. $K$ 的非零元素是 $K$ 中的 $n$ 次幂, 只要它对除可能有限多个之外的每个素 $v$ 都是 $K_v$ 中的 $n$ 次幂.
\end{theorem}

\begin{proof}
回顾: 对包含本原 $n$ 次单位根的域 $k$, 多项式 $X^{n}-a$ 只要在 $k$ 中有一个根, 就在 $k[X]$ 中完全分裂. 设 $a$ 是 $K$ 的非零元素, 并设 $\beta$ 是 $a$ 在某个扩域中的 $n$ 次根. 若 $a$ 对几乎所有 $v$ 都是 $K_v$ 中的 $n$ 次幂, 则 $X^{n}-a$ 对几乎所有 $v$ 在 $K_v[X]$ 中完全分裂, 从而 $v$ 在 $K[\beta]$ 中完全分裂. 由于这对几乎所有 $v$ 成立, (VII, \ref{prop:vii.4.6}) 表明 $K[\beta]=K$.
\end{proof}

\begin{remark}\label{rem:viii.1.2}
如果不用 (VII, \ref{prop:vii.4.6}) 而用 (VI, \ref{thm:vi.3.4}), 则得到更强的结果: 在该定理的假设下, $K$ 的非零元素是 $K$ 中的 $n$ 次幂, 只要它在密度 $>1/2$ 的一组素 $v$ 处成为 $K_v$ 中的 $n$ 次幂.
\end{remark}

人们或许希望该定理在 $K$ 不含本原 $n$ 次单位根时仍然成立, 但下述习题表明并非如此, 甚至对 $K=\Q$ 与 $n=8$ 都不成立.

\begin{exercise}\label{exc:viii.1.3}
证明 $16$ 在 $\R$ 与所有奇素数 $p$ 的 $\Q_p$ 中都是 $8$ 次幂, 但 (古希腊人早已知道) 它在 $\Q$ 中不是. (提示: 证明 $\Q[\zeta_8]=\Q[i,\sqrt{2}]$ 在所有奇素数处不分歧, 并由此推出: 对奇数 $p$, $\Q_p$ 至少包含 $1+i$, $\sqrt{2}$, $\sqrt{-2}$ 三者之一.)
\end{exercise}

域 $\Q[\zeta_8]$ 在 $\Q$ 上不是循环的, 故该习题与下述定理并不矛盾.

\begin{theorem}\label{thm:viii.1.4}
设 $K$ 是数域, 并设 $n$ 为整数, 使得 $K[\zeta_{2^{t}}]$ 在 $K$ 上是循环的, 其中 $2^{t}$ 是整除 $n$ 的 2 的幂. $K$ 的非零元素是 $K$ 中的 $n$ 次幂, 只要它对除可能有限多个之外的每个素 $v$ 都是 $K_v$ 中的 $n$ 次幂.
\end{theorem}

\begin{proof}
第一步. 不妨设 $n$ 为素数幂来证明该定理. 设 $n=n_1n_2$, 其中 $n_1$ 与 $n_2$ 互素, 并设定理对 $n_1$ 与 $n_2$ 成立. 于是对几乎所有 $v$ 都是 $K_v$ 中 $n$ 次幂的非零元素 $a\in K$ 既是 $K$ 中的 $n_1$ 次幂又是 $K$ 中的 $n_2$ 次幂, 比如 $a=b^{n_1}$, $a=c^{n_2}$. 现在存在整数 $r,s$ 使得 $rn_1+sn_2=1$, 故
\[
a=a^{rn_1}a^{sn_2}=c^{rn_1n_2}b^{sn_1n_2}\in K^{\times n}.
\]

第二步. 当 $n$ 是素数 $p$ 的幂且 $K[\zeta_n]/K$ 是 $p$ 幂阶的循环扩张时, 定理成立. 令 $K'=K[\zeta_n]$, 并设 $a$ 是对几乎所有 $v$ 成为 $K_v$ 中 $n$ 次幂的非零元素. 根据定理\ref{thm:viii.1.1}, $a$ 在 $K'$ 中成为 $n$ 次幂, 比如 $a=\beta^{n}$, 于是
\[
X^{n}-a=\prod_{j=0}^{n-1}(X-\beta\zeta_n^{j})\qquad\text{在 } K'[X] \text{ 中}.
\]
设
\[
X^{n}-a=\prod_{i}f_{i}(X)
\]
是 $X^{n}-a$ 在 $K[X]$ 中分解为不可约因子的分解. 对每个 $i$, 取 $f_i(X)$ 的一个根 $\beta_i=\beta\zeta_n^{j(i)}$. 则 $K[\beta_i]\subset K'$, 从而 $K[\beta_i]$ 是 $K$ 的 Abel 扩张, 特别地, 它在 $K$ 上是 Galois 的, 因而是 $f_i(X)$ 的分裂域. 设 $v$ 是使 $a$ 为 $K_v$ 中 $n$ 次幂的 $K$ 的素. 由假设, $X^{n}-a$ 在 $K_v$ 中有根, 故至少某个 $f_i(X)$ 在 $K_v$ 中有根. 对那个特定的 $i$, $v$ 在 $K[\beta_i]$ 中完全分裂. 于是我们看到: 每个 $v\notin S$ 至少在诸域 $K[\beta_i]$ 之一中完全分裂, 但不同的 $v$ 可能在不同的域中分裂, 由此我们得不出任何结论. 要看出所有的 $v$ 都在同一个 $K[\beta_i]$ 中分裂, 必须用到 $K'$ 在 $K$ 上是素数幂阶循环的这一假设. 该假设蕴含中间域是线性有序的:
\[
K'\supset\cdots\supset K_{3}\supset K_{2}\supset K_{1}\supset K.
\]
取 $i_0$ 使 $K[\beta_{i_0}]$ 是诸 $K[\beta_i]$ 中最小的一个. 则每个 $v\notin S$ 在一个包含 $K[\beta_{i_0}]$ 的域中完全分裂, 从而在 $K[\beta_{i_0}]$ 中完全分裂. 故 $K[\beta_{i_0}]=K$, 即 $X^{n}-a$ 在 $K$ 中至少有一个根.

第三步. 定理成立. 经第一步, 可设 $n=p^{r}$; 经第二步, 可设 $p$ 为奇数. 设 $a$ 对几乎所有 $v$ 是 $K_v$ 中的 $n$ 次幂. 由于 $K[\zeta_{p^r}]$ 在 $K[\zeta_p]$ 上是 $p$ 幂阶的循环扩张, 第二步表明 $a$ 在 $K[\zeta_p]$ 中成为 $n$ 次幂, 比如 $a=b^{n}$. 取范数, 得知 $a^{d}$ 是 $K^{\times}$ 中的 $n$ 次幂, 其中 $d=[K[\zeta_p]:K]<p$. 但 $d$ 与 $p$ 互素, 故这蕴含 $a$ 是 $K^{\times}$ 中的 $n$ 次幂 (我们知道 $a^{d}=1$ 属于 $K^{\times}/K^{\times p^{r}}$, 而这蕴含 $a=1$ 属于 $K^{\times}/K^{\times p^{r}}$).
\end{proof}

\begin{remark}\label{rem:viii.1.5}
(a) 对 $K$ 的素的有限集 $S$ 与整数 $n$, 令 $P(n,S)$ 为所有满足如下条件的 $a\in K^{\times}$ 构成的群: 对所有 $v\notin S$, $a\in K_v^{\times n}$. 对 $2^{t}$ 情形作进一步分析可以发现: 要么 $P(n,S)=K^{\times n}$, 要么存在 $a_0\in K^{\times}$ 使 $P(n,S)=K^{\times n}\sqcup a_0K^{\times n}$. 例如, 对 $K=\Q$, $P(8,S)=\Q^{\times 8}\sqcup 16\,\Q^{\times 8}$ (\cite{ArtinTate1961}, p.~96; 2009 年版 p.~75). 注意 $P(n,S)/K^{\times n}$ 是使
\[
0\to P(n,S)/K^{\times n}\to K^{\times}/K^{\times n}\to\mathbb{I}_{K}/\mathbb{I}_{K}^{n}\to\mathbf{C}_{K}/\mathbf{C}_{K}^{n}\to0
\]
正合的群.

(b) 该定理的表述并非对所有 $K,S,n$ 都成立, 这一事实可视为小数域的一种病态. 例如, 若 $\sqrt{-1}\in K$, 则 $K[\zeta_{2^t}]$ 在 $K$ 上是循环的.

(c) 证明中的关键一步是第二步, 它有点微妙. Whaples 在证明 Grunwald ``定理''时正是在这里出错的: 他用 $a^{\frac{1}{n}}$ 表示 $X^{n}-a$ 的一个根, 却忘了不同的根相对于 $K$ 可能有十分不同的性质 (例如, 它们的极小多项式是诸 $f_i$, 甚至可能有不同的次数). Artin 与 Tate 明确警告不要试图缩短上述论证.
\end{remark}

\section{Grunwald--Wang 定理}\label{sec:VIII.2}

类域论通过把关于 Abel 扩张的问题转化为关于 idèle 类群中有限指数开子群的问题, ``解决''了所有这类问题. 遗憾的是, 关于 idèle 类群中有限指数开子群的问题也可能很难. 本节考虑如下问题.

\begin{question}\label{qu:viii.2.1}
设 $K$ 是代数数域, 并设对有限集 $S$ 中的每个素 $v$ 给定了 $K_v$ 的 $n_v$ 次有限循环扩张 $K^{v}$. 是否存在 $K$ 的 $n=\mathrm{lcm}(n_v)$ 次有限循环扩张 $L/K$, 使得 $L^{v}\approx K^{v}$? (此条件是指: 存在 $L$ 的整除 $v$ 的素 $w$ 以及 $K_v$-同构 $L_w\to K^{v}$.)
\end{question}

下面的例子表明答案并不总是肯定的, 然而我们将会看到, 答案通常是肯定的.

\begin{example}\label{ex:viii.2.2}
设 $\Q^{2}$ 是 $\Q_2$ 的 $8$ 次非分歧扩张. 容易构造 $\Q$ 的 $8$ 次扩张 $L$ 使得 $(2)$ 在 $L$ 中保持素, 从而对唯一的 $w|2$ 有 $L_w\approx\Q^{2}$: 取 $L=\Q[X]/(g(X))$, 其中 $g(X)$ 是 $\Z[X]$ 中任意模 $2$ 不可约的首一 $8$ 次多项式. 然而, 我断言不可能构造出这样的循环扩张. 更精确地说:

设 $L$ 是 $\Q$ 的 $8$ 次循环扩张, 则 $(2)$ 不可能在 $L$ 中保持素.

这由如下三个断言推出:

(i) $16$ 在所有 $p\ne 2$ 的 $\Q_p$ 中都是 $8$ 次幂 (包括 $p=\infty$), 但在 $\Q_2$ 中不是.

(ii) 设 $L/\Q$ 是有限 Abel 扩张; 若 $a\in\Q$ 对除一个之外的所有 $p$ 都成为 $\Q_p$ 中的范数, 则它对所有 $p$ 都是范数.

(iii) $16$ 不是来自 $\Q^{2}$ 的范数.

设 $L$ 是 $\Q$ 的 $8$ 次循环扩张. 由于 $L_v$ 在 $\Q_p$ 上的次数整除 $8$, (i) 表明 $16$ 对所有 $p\ne 2$ 都成为 $\Q_p$ 中的范数. 根据 (ii), 这蕴含 $16$ 是 $\Q_2$ 中的范数. 现在 (iii) 表明 $L^{2}$ 不同构于 $\Q^{2}$.

关于断言 (i) 的证明, 见习题\ref{exc:viii.1.3} (及其解答). 对 (ii), 考虑图表:
\begin{equation*}
\begin{tikzcd}[column sep=small, row sep=scriptsize]
& \Q_p^{\times} \ar[dr, "\phi_p"] & \\
\Q^{\times} \ar[ur, hook] \ar[r] & \mathbb{I}_{\Q} \ar[r, "\phi"] \ar[d, two heads] & \Gal(L^{p}/\Q_p) \\
& \Q_p^{\times} \ar[ur, two heads] &
\end{tikzcd}
\end{equation*}
其中 $\phi_p$ 是到整除 $p$ 的素处的分解群中的局部 Artin 映射, $\phi$ 是整体 Artin 映射 (我们允许 $p=\infty$, 此时 $\Q_p=\R$). 图表交换 (这正是我们定义整体 Artin 映射的方式). 对每个 $a\in\Q^{\times}$, 有乘积公式
\[
\prod_{p=2,3,5,\ldots,\infty}\phi_{p}(a)=1
\]
(VII, \ref{cor:vii.4.8}). 现在 $\phi_p(a)=1$ 当且仅当 $a$ 成为 $\Q_p$ 中的范数, 故乘积公式蕴含 (ii). 对 (iii), 我们利用: 由于 $\Q^{2}$ 在 $\Q_2$ 上非分歧, $\Q_2$ 上的正规赋值延拓为 $\Q^{2}$ 上的正规赋值. 若 $16=\Nm(\alpha)$, 则
\[
4=\mathrm{ord}_{2}(16)=\mathrm{ord}_{2}(\mathrm{Nm}(\alpha))=8\,\mathrm{ord}_{2}(\alpha),
\]
这不可能.
\end{example}

我现在简要解释一下正面的结果 (\cite{ArtinTate1961}, 第 9、10 章). 设 $S$ 是 $K$ 的素的有限集. 复合映射
\[
\prod_{v\in S}K_{v}^{\times}\to\mathbb{I}_{K}\to\mathbf{C}_{K}
\]
是单射, 并且当我们赋予 $P\overset{\text{def}}{=}\prod_{v\in S}K_v^{\times}$ 积拓扑时它是连续的. 当 $S$ 含有多于一个素时, 积拓扑与由 $\mathbf{C}_K$ 诱导的拓扑并不一致 (同上, p.~98). 尽管如此, $P$ 的一个有限指数子群对一种拓扑是开的当且仅当它对另一种拓扑是开的, 并且 $P$ 的有限指数开子群是 $P$ 与 $\mathbf{C}_K$ 中某个有限指数开子群的交 (同上, p.~99). 利用互反律, 这蕴含下述结果.

\begin{theorem}\label{thm:viii.2.3}
对每个 $v\in S$, 设 $N_v$ 是 $K_v^{\times}$ 中有限指数的开子群. 则存在 $K$ 的 Abel 扩张 $L$, 使得对所有 $v\in S$, $L^{v}$ 是 (由局部类域论) 对应于 $N_v$ 的 $K_v$ 的扩张.
\end{theorem}

更细致的分析证明下述结果 (其中 ``特征''指``连续特征'').

\begin{theorem}[GRUNWALD--WANG]\label{thm:viii.2.4}
对每个 $v\in S$, 设 $\chi_v$ 是 $K_v^{\times}$ 的特征. 则存在 $\mathbf{C}_K$ 的特征 $\chi$, 它在 $K_v^{\times}$ 上的限制为 $\chi_v$ 对所有 $v\in S$ 成立. 设 $n_v$ 为 $\chi_v$ 的阶, $n=\mathrm{lcm}(n_v)$. 则可以选取 $\chi$ 使其阶为 $n$, 但当 $2^{t}|n$ 对某个使 $K[\zeta_{2^{t}}]$ 在 $K$ 上不循环的 $t$ 成立时可能例外.
\end{theorem}

\begin{proof}
见 \cite{ArtinTate1961}, 第 X 章, 定理 5. 或者, 也可以通过比较 $S$ 与全体素集的 Poitou--Tate 定理来证明该定理的一个更精确的版本——见这些讲义的未来版本.
\end{proof}

例\ref{ex:viii.2.2} 表明, 若去掉对 $n$ 的条件, 第二个断言可能失效.

\begin{corollary}\label{cor:viii.2.5}
对每个 $v\in S$, 设 $n_v$ 为正整数. 若 $v$ 是无限素, 则 $n_v$ 必须是 $K_v$ 的某个扩张的可能次数. 则存在 $K$ 的 $n=\mathrm{lcm}(n_v)$ 次循环扩张 $L$, 使得局部次数对所有 $v\in S$ 为 $n_v$.
\end{corollary}

\begin{proof}
对无限素 $v\in S$, 以唯一可能的方式选取 $K_v^{\times}$ 上 $n_v$ 阶的特征 $\chi_v$; 对不整除 $2$ 的有限素 $v\in S$, 取 $\chi_v$ 为在 $U_v$ 上取值 $1$ 的 $n_v$ 阶典范特征; 对整除 $2$ 的 $v\in S$, 选取 $n_v$ 阶的 $\chi_v$ 以避开定理中的例外情形. 见同上 p.~105, 定理 5 (不要与\footnote{Serge Lang 记录了本讲义. 在 2009 年版中, 它已被正确编号为定理 6.}同页码 p.~103 上编号相同的定理混淆).
\end{proof}

\paragraph{应用}
设 $D$ 是代数数域 $K$ 上次数为 $d^{2}$ 的中心可除代数, 并设 $i_v$ 是 $D\otimes_K K_v$ 的类在不变映射 $\mathrm{inv}_v:\mathrm{Br}(K_v)\to\Q/\Z$ 下的像. 写 $i_v=m_v/n_v \bmod \Z$, 其中 $m_v,n_v$ 为互素整数且 $n_v>0$.

由于 $\mathrm{Br}(K_v)\to\Q/\Z$ 是单射 (当 $v$ 有限时甚至是同构), $D\otimes_K K_v$ 的类在 $\mathrm{Br}(K_v)$ 中的阶为 $n_v$. 在同构意义下, 恰好存在一个 $K_v$ 上不变量为 $i_v=m_v/n_v$ 的可除代数 $D_{i_v}$, 其次数为 $n_v^{2}$. 于是 $D\otimes_K K_v\approx M_{d/n_v}(D_{n_v})$. 特别地, $n_v|d$, 从而 $n|d$.

由于 $\mathrm{Br}(K)\to\bigoplus_v\mathrm{Br}(K_v)$ (对所有 $v$ 求和) 是单射, $D$ 的类在 $\mathrm{Br}(K)$ 中的阶为 $n=\mathrm{lcm}(n_v)$.

根据推论, 存在 $K$ 的 $n$ 次循环扩张 $L$, 其局部次数对每个满足 $i_v\notin\Z$ 的 $v$ 为 $n_v$. 于是 $L$ 分裂 $D$. 由 (IV, \ref{cor:iv.3.6}), 存在 $K$ 上的中心单代数 $B$ 使得

(a) $B\supset L$;

(b) $B$ 相似于 $D$, 即 $B\approx M_{r}(D)$ 对某个 $r$; 且

(c) $[B:K]=[L:K]^{2}$.

由后两条可见 $r^{2}d^{2}=n^{2}$. 由于 $n|d$, 故实际上 $r=1$, $n=d$. 我们证明了:

\begin{theorem}\label{thm:viii.2.6}
对数域 $K$ 上的每个可除代数 $D$, $D$ 在 $\mathrm{Br}(K)$ 中的阶为 $\sqrt{[D:K]}$, 并且 $D$ 包含 $K$ 的一个循环扩张作为极大子域 (从而被它分裂).
\end{theorem}

由该定理与 Noether--Skolem 定理可以推出: $D$ 是 Dickson 于 1906 年所定义意义下的循环代数, 即 $D$ 作为 $K$-代数由 $L$ 的元素与单个元素 $\beta$ 生成, 满足关系
\[
\beta\cdot\alpha=\sigma\alpha\cdot\beta,\quad \alpha\in L,\qquad \beta^{d}=\gamma,
\]
其中 $\sigma$ 生成 $\Gal(L/K)$, $\gamma\in K^{\times}$.

这些是非常著名的结果\footnote{``1930 年代代数结构理论的最高成就, 无疑是如下定理: 数域上每个有限维中心可除代数都是循环的, 以及用一组数值不变量对这些代数的分类. 后一结果等价于确定数域 Brauer 群的结构. 除了我们已概述的中心单代数的一般理论之外, 这些基本结果的证明还需要 Hasse 关于 $p$-adic 域上中心单代数的结构理论、Hasse 范数定理 (``Hasse 原理''), 以及关于某些循环扩张的 Grünwald 存在定理……数域上中心可除代数循环性的第一个证明由 Brauer、Hasse 与 Noether 给出 (1931).'' Nathan Jacobson, A.A.~Albert 讣告, \emph{Bull.~AMS}, 80 (1974), p.~1081.}, 有许多应用, 例如在 Abel 簇理论中的应用 (见我的讲义 AV, 定理 16.8).

\begin{NOTES}
1933 年, Grunwald 在不带对 $n$ 的条件的情况下证明了定理\ref{thm:viii.2.4}; 1942 年, Whaples 给出了第二个证明. 1948 年, Wang 发现了由 (\ref{exc:viii.1.3}, \ref{ex:viii.2.2}) 提供的反例, 并修正了表述. 在 1933 年之后的十五年间, Grunwald 的``定理''被相当广泛地使用, 例如用来证明数域上的每个可除代数都是循环的. 见 \cite{Roquette2005}, 5.3.
\end{NOTES}

事实上, 正如 Brian Conrad 向我指出的, Grunwald 定理的一个反例实际上早在 1934 年, 即比 Wang 早 14 年就已发表, 但 (显然) 没有人——包括作者本人——注意到这一点. 具体地说, 设 $a,n$ 为整数; 对不整除 $an$ 的素数 $p$, $a$ 是 $\Q_p$ 中的 $n$ 次幂当且仅当它模 $p$ 是 $n$ 次幂. Trost (1934)\footnote{E.~Trost, Zur Theorie der Potenzreste, Nieuw Arch. Wisk. 18 (1934), 58--61.}用初等方法证明了下述定理:

若 $a$ 对几乎所有素数都是 $n$ 次幂且 $8$ 不整除 $n$, 则 $a$ 是 $n$ 次幂; 当 $8$ 整除 $n$ 时, $a$ 要么是 $n$ 次幂, 要么是 $2^{n/2}$ 乘以一个 $n$ 次幂 (两种情形都会出现).

例如, 当 $n=8$ 时, 有可能 $2^{4}a$ 是 $8$ 次幂 (从而 $a$ 不是 $8$ 次幂). 这在 $a=16$ 时发生. 换言之, Trost 已经知道例\ref{ex:viii.2.2} 的断言 (i), 而由它通过标准论证即可导出 Grunwald 定理的反例.

顺带一提, Olga Taussky 审阅过 Grunwald 的论文. Chevalley 为 Mathematical Reviews 审阅过 Whaples 的论文, 甚至仔细到注意到``在阅读基本引理的证明时须提防若干排版错误'', 却没有仔细到发现主要结果是错误的.
\section{范数与二次型的局部--整体原则}\label{sec:VIII.3}

局部--整体 (或 Hasse) 原则问的是: 一个命题若在数域 $K$ 的每个完备化上都成立, 它是否就在 $K$ 上成立. 本节给出类域论使我们能够证明该原则成立的两个情形.

\subsection*{范数}

\begin{theorem}[HASSE 范数定理]\label{thm:viii.3.1}
设 $L/K$ 是数域的循环扩张, $a\in K^{\times}$. 则 $a$ 在 $K_v$ 中的像对除有限多个之外的所有 $v$ 都是来自 $L^{v}$ 的范数; 并且若它对所有 $v$ 都是范数, 则它是 $K$ 中的范数.
\end{theorem}

\begin{proof}
根据定理 (VII, \ref{thm:vii.5.1}), $H^{1}(G,\mathbf{C}_{L})=0$. 由于循环群上同调的周期性, 这蕴含 $H_{T}^{-1}(G,\mathbf{C}_{L})=0$. 因此, 从
\[
1\to L^{\times}\to\mathbb{I}_{L}\to\mathbf{C}_{L}\to0
\]
的上同调列得到
\[
H_{T}^{0}(G,L^{\times})\to H_{T}^{0}(G,\mathbb{I}_{L})
\]
是单射. 但 (见 (VII, \ref{prop:vii.2.5})) 这正是映射
\[
K^{\times}/\Nm(L^{\times})\to\bigoplus_{v}K_{v}^{\times}/\Nm(L^{v\times}).
\]
\end{proof}

\begin{remark}\label{rem:viii.3.2}
该证明对非循环扩张失效, 并且事实上该命题对非循环扩张并不成立. 例如, $5^{2}$ 在所有素处都是来自 $\Q[\sqrt{13},\sqrt{17}]$ 的局部范数, 但不是整体范数 (见 \cite{CasselsFrohlich1967}, 习题 5.3, p.~360). 对给定的非循环 Abel 群 $G$, 下述论文研究以 $G$ 为 Galois 群且满足 Hasse 范数原理的扩张的密度: Frei, C.; Loughran, D.; Newton, R.: The Hasse norm principle for Abelian extensions. \emph{Amer.~J.~Math.} 140 (2018), no.~6, 1639--1685.
\end{remark}

\subsection*{二次型}

回顾: 域 $k$ 上向量空间 $V$ 上的二次型是指映射 $Q\colon V\to k$ 使得

(a) $Q(av)=a^{2}Q(v)$;

(b) $B(v,w)\overset{\text{def}}{=}Q(v+w)-Q(v)-Q(w)$ 是 $V$ 上的双线性型.\footnote{当 $k$ 的特征 $\ne 2$ 时, 通常令 $B(v,w)=\frac{1}{2}(Q(v+w)-Q(v)-Q(w))$.}

二次型 $Q$ 称为\textbf{非退化的} (nondegenerate), 若其伴随双线性型 $B$ 非退化. 设 $c\in k$. 非退化二次型 $Q$ 称为\textbf{表示} (represent) $c$, 若存在非零向量 $v\in V$ 使 $Q(v)=c$.

\begin{lemma}\label{lem:viii.3.3}
若非退化二次型 $Q$ 表示 $0$, 则它表示所有 $c\in k$.
\end{lemma}

\begin{proof}
注意, 对 $t\in k$,
\[
Q(tv+w)=t^{2}Q(v)+tB(v,w)+Q(w).
\]
设 $v_0$ 是满足 $Q(v_0)=0$ 的非零向量, 则因 $B$ 非退化, 存在向量 $w_0$ 使 $B(v_0,w_0)\ne 0$. 当 $t$ 取遍 $k$ 的所有值时, $Q(tv_0+w_0)=tB(v_0,w_0)+Q(w_0)$ 也取遍 $k$ 的所有值.
\end{proof}

当 $k$ 的特征 $\ne 2$ 时, 存在 $V$ 的一组基 $\{e_1,\ldots,e_n\}$ 使 $B(e_i,e_j)=0$ 对 $i\ne j$ 成立. 此时
\[
Q\Big(\sum x_ie_i\Big)=\sum_{i=1}^{n}a_ix_i^{2},\qquad n=\dim V.
\]
以下我们将二次型写作
\[
q(X_{1},\ldots,X_{n})=a_{1}X_{1}^{2}+\cdots+a_{n}X_{n}^{2},
\]
并记住: 可逆的变量代换不会改变我们的任何结论.

\begin{lemma}\label{lem:viii.3.4}
非退化二次型 $q(X_{1},\ldots,X_{n})$ 表示 $a$ 当且仅当 $r\overset{\text{def}}{=}q-aY^{2}$ 表示 $0$.
\end{lemma}

\begin{proof}
若 $q(x_1,\ldots,x_n)=a$, 则 $r(x_1,\ldots,x_n,1)=0$. 反之, 设 $r(x_1,\ldots,x_n,y)=0$. 若 $y=0$, 则 $q$ 表示 $0$, 从而表示 $k$ 的每个元素. 若 $y\ne 0$, 则 $q\big(\frac{x_1}{y},\ldots,\frac{x_n}{y}\big)=q(x_1,\ldots,x_n)/y^{2}=a$.
\end{proof}

\begin{theorem}\label{thm:viii.3.5}
设 $q$ 是系数在数域 $K$ 中的 $n$ 元非退化二次型.

(a) 若 $n\ge 3$, 则 $q$ 对除有限多个之外的所有 $v$ 在 $K_v$ 中表示 $0$.

(b) $q$ 在 $K$ 中表示 $0$ 当且仅当它对所有 $v$ 在 $K_v$ 中表示 $0$.
\end{theorem}

在开始证明之前, 我们指出一个推论.

\begin{corollary}\label{cor:viii.3.6}
设 $c\in K$. 系数在 $K$ 中的非退化二次型 $q$ 在 $K$ 中表示 $c$ 当且仅当它对所有 $v$ 在 $K_v$ 中表示 $c$.
\end{corollary}

\begin{proof}
令 $r=q-cY^{2}$. 则 $r$ 表示 $0$ 当且仅当 $q$ 表示 $c$.
\end{proof}

我们从一个纯代数的结果开始证明.

\begin{proposition}\label{prop:viii.3.7}
设 $k$ 是特征 $\ne 2$ 的域.

(a) 型 $q=X^{2}$ 不表示 $0$.

(b) 型 $q=X^{2}-aY^{2}$ 表示 $0$ 当且仅当 $a$ 是平方.

(c) 型 $q=X^{2}-aY^{2}-bZ^{2}$ 表示 $0$ 当且仅当 $b$ 是来自域 $k[\sqrt{a}]$ 的范数.

(d) 型 $q=X^{2}-bY^{2}-cZ^{2}+acT^{2}$ 在 $k$ 中表示 $0$ 当且仅当 $c$ 作为 $k[\sqrt{ab}]$ 的元素是来自 $k[\sqrt{a},\sqrt{b}]$ 的范数.
\end{proposition}

\begin{proof}
(a) 显然.

(b) 根据引理\ref{lem:viii.3.4}, $X^{2}-aY^{2}$ 表示 $0$ 当且仅当 $X^{2}$ 表示 $a$.

(c) 根据引理\ref{lem:viii.3.4}, $X^{2}-aY^{2}-bZ^{2}$ 表示 $0$ 当且仅当 $X^{2}-aY^{2}$ 表示 $b$, 即 $b$ 是来自 $k[\sqrt{a}]$ 的范数.

(d) 显然,
\[
q(x,y,z,t)=0\iff c=\frac{\Nm_{k[\sqrt{b}]/k}(x+\sqrt{b}\,y)}{\Nm_{k[\sqrt{a}]/k}(z+\sqrt{a}\,t)}.
\]
由于范数的逆仍是范数, 这表明 $q$ 表示零当且仅当 $c$ 是一个来自 $k[\sqrt{a}]$ 的范数与一个来自 $k[\sqrt{b}]$ 的范数之积. 于是 (d) 由下述引理推出.
\end{proof}

\begin{lemma}\label{lem:viii.3.8}
设 $k$ 是特征 $\ne 2$ 的域. 元素 $c\in k^{\times}$ 是一个来自 $k[\sqrt{a}]$ 的范数与一个来自 $k[\sqrt{b}]$ 的范数之积, 当且仅当它作为 $k[\sqrt{ab}]$ 的元素是来自 $L=k[\sqrt{a},\sqrt{b}]$ 的范数.
\end{lemma}

\begin{proof}
我们把 $a$、$b$ 或 $ab$ 是 $k$ 中平方的退化情形留给读者. 于是可设 $\Gal(k[\sqrt{a},\sqrt{b}]/k)=\{1,\sigma,\tau,\sigma\tau\}$, 其中 $\sigma$、$\tau$、$\sigma\tau$ 的阶都是 $2$, 并分别固定 $\sqrt{a}$、$\sqrt{b}$、$\sqrt{ab}$. 第一个条件断言:

(*) 存在 $x,y\in k[\sqrt{a},\sqrt{b}]$ 使得 $\sigma x=x$, $\tau y=y$, 且 $xy\cdot\sigma\tau(xy)=c$;

第二个条件断言:

(**) 存在 $z\in k[\sqrt{a},\sqrt{b}]$ 使得 $z\cdot\sigma\tau(z)=c$.

显然 (*) $\Rightarrow$ (**). 反之, 注意
\[
z\cdot\sigma z=\Nm_{k[\sqrt{a},\sqrt{b}]/k[\sqrt{a}]}z\in k[\sqrt{a}].
\]
再者,
\[
\Nm_{k[\sqrt{a}]/k}(z\cdot\sigma z)=z\cdot\sigma z\cdot\tau z\cdot\sigma\tau z\in k.
\]
由于 $z\cdot\sigma\tau z=c\in k$, 这蕴含 $\sigma z\cdot\tau z\in k$, 故
\[
\sigma z\cdot\tau z=\sigma(\sigma z\cdot\tau z)=z\cdot\sigma\tau z=c.
\]
因此
\[
\Nm_{k[\sqrt{a}]/k}(z\cdot\sigma z)=c^{2}.
\]
现在把 Hilbert 定理 90 ((II, \ref{cor:ii.1.23})) 应用于 $z\cdot\sigma z/c\in k[\sqrt{a}]$, 得知存在 $x\in k[\sqrt{a}]^{\times}$ 使 $\tau x/x=z\cdot\sigma z/c$. 令 $y=\sigma\tau z/x$. 则
\[
\tau y=\sigma z/\tau x=c/(z\cdot x)=z\cdot\sigma\tau z/(z\cdot x)=y
\]
(依次利用: $y$ 的定义; $x$ 的定义; $z$ 的定义; $y$ 的定义), 并且
\[
xy\cdot\sigma\tau(xy)=\sigma\tau z\cdot\sigma\tau(\sigma\tau z)=\sigma\tau z\cdot z=c
\]
(依次利用: $y$ 的定义; $(\sigma\tau)^{2}=1$; $z$ 的定义), 如所欲证.
\end{proof}

\begin{proof}[定理\ref{thm:viii.3.5}(a) 的证明]
若 $q=q_1(X_1,\ldots,X_m)+q_2(X_{m+1},\ldots,X_n)$ 且 $q_1$ 表示零, 则 $q$ 也表示零. 因此只需对三元二次型证明 (\ref{thm:viii.3.5}a). 用非零纯量乘 $q$ 之后, 可设 $q=X^{2}-aY^{2}-bZ^{2}$, 而对此种二次型, 结论由 (\ref{prop:viii.3.7}) 与定理 \ref{thm:viii.3.1} 推出.
\end{proof}

\begin{proof}[定理\ref{thm:viii.3.5}(b) 的证明]
我们对变量个数 $n$ 作归纳证明.

当 $n=1$ 时, 无需证明, 因为假设永远不会成立.

当 $n=2$ 时, 用非零纯量乘 $q$ 之后可设 $q=X^{2}-aY^{2}$, 对此种二次型, 结论由 (\ref{prop:viii.3.7}) 与定理\ref{thm:viii.1.1} 推出.

当 $n=3,4$ 时, 结论以类似方式由 (\ref{prop:viii.3.7}c,d) 与定理\ref{thm:viii.1.1} 推出.

在证明一般情形之前, 我们作一些初等观察.

(a) 非退化二次型 $q_1(X_1,\ldots,X_m)-q_2(X_{m+1},\ldots,X_n)$ 在域 $k$ 中表示 $0$ 当且仅当存在 $c\in k$ 使 $q_1$ 与 $q_2$ 都表示 $c$.

(b) 若 $q$ 在 $k^\times$ 中表示 $c$, 则 $q$ 表示傍系 $c\cdot k^{\times 2}$ 中的每个元素.

(c) 子群 $K_v^{\times 2}\subset K_v^{\times}$ 是开的. 当 $v$ 为实或复时显然. 当 $v$ 非阿基米德时, Newton 引理 (ANT, 7.32) 表明, 对任意满足 $|1-a|_{v}<|2|_{v}^{2}$ 的 $a$, $1$ 可以精化为 $X^{2}-a$ 的根.

结合 (b) 与 (c) 可见: 系数在 $K_v$ 中的二次型 $q$ 表示 $K_v^{\times}$ 的一个非空开子集中的元素.

现在设 $n\ge 5$, 并设定理\ref{thm:viii.3.5}(b) 已对 $n-1$ 获证. 令
\[
q(X_{1},\ldots,X_{n})=aX_{1}^{2}+bX_{2}^{2}-r(X_{3},\ldots,X_{n}),\qquad n-2\ge 3.
\]
由定理的 (a) 知: 除某个有限集 $S$ 中的 $v$ 之外, $r$ 在 $K_v$ 中表示 $0$. 设 $v\in S$. 由于 $q$ 在 $K_v$ 中表示 $0$, 存在同时被 $aX_1^{2}+bX_2^{2}$ 与 $r$ 表示的元素 $c_v\in K_v^{\times}$, 即存在 $x_i(v)\in K_v$ 使
\[
ax_{1}(v)^{2}+bx_{2}(v)^{2}=c_v=r(x_{3}(v),\ldots,x_{n}(v)).
\]
现在应用弱逼近定理 (ANT, 7.20), 找到对所有 $v\in S$ 都在 $K_v$ 中接近 $x_1(v)$、$x_2(v)$ 的元素 $x_1,x_2\in K$. 则
\[
c\overset{\text{def}}{=}ax_{1}^{2}+bx_{2}^{2}
\]
将对每个 $v\in S$ 接近 $c_v$; 特别地, 可设 $c/c_v\in K_v^{\times 2}$ 对所有 $v\in S$ 成立.

考虑二次型 $cY^{2}-r$. 它在 $K_v$ 中表示零, 对 $v\notin S$: 因为 $r$ 在 $K_v$ 中表示零; 对 $v\in S$: 因为 $r$ 在 $K_v$ 中表示 $c$. 由归纳, $cY^{2}-r$ 在 $K$ 中表示零. 由此得到 $q$ 在 $K$ 中表示 $0$, 因为 $aX_1^{2}+bX_2^{2}$ 与 $r$ 都在 $K$ 中表示 $c$.
\end{proof}

\begin{proposition}\label{prop:viii.3.9}
$\Q_p$ 的有限扩张 $K$ 上的四元非退化二次型 $q$ 表示 $K^{\times}$ 的每个非零元素.
\end{proposition}

\begin{proof}
若 $q$ 表示 $0$, 则它表示 $K$ 的每个元素. 我们假设相反的情形. 用 $K$ 的非零元素乘 $q$ 之后, 可设
\[
q=X^{2}-bY^{2}-cZ^{2}+acT^{2}.
\]
由于 $q$ 在 $K$ 中不表示 $0$, $b$ 与 $a$ 都不是平方.

若 $K[\sqrt{a}]\ne K[\sqrt{b}]$, 则 (由局部类域论定理 (I, \ref{thm:1.1})) $\Nm(K[\sqrt{a}]^{\times})$ 与 $\Nm(K[\sqrt{b}]^{\times})$ 是 $K^{\times}$ 中两个不同的指数 $2$ 子群, 从而 $K^{\times}=\Nm(K[\sqrt{a}]^{\times})\cdot\Nm(K[\sqrt{b}]^{\times})$. 由于范数的逆仍是范数, 这意味着可以把 $c$ 写成
\[
c=\frac{x^{2}-by^{2}}{z^{2}-at^{2}},
\]
其中 $x,y,z,t\in K$. 展开相乘后发现 $q$ 表示 $0$, 与假设矛盾. 因此 $K[\sqrt{a}]=K[\sqrt{b}]$, 且 $a=b\times(\text{平方})$ (见 (VII, \ref{prop:vii.A.1})). 该平方可被吸收进 $T^{2}$, 故可写
\[
q=X^{2}-bY^{2}-cZ^{2}+bcT^{2}.
\]
考虑四元数代数 $H(b,c)$ (见习题 IV.\ref{ex:iv.5.7}). 对
\[
\alpha=x+yi+zj+tk
\]
定义
\[
\bar{\alpha}=x-yi-zj-tk,
\]
于是
\[
\Nm(\alpha)\overset{\text{def}}{=}\alpha\bar{\alpha}=x^{2}-by^{2}-cz^{2}+bct^{2}.
\]
映射 $\alpha\mapsto\Nm(\alpha):H(b,c)^{\times}\to K^{\times}$ 是同态, 我们必须证明它是满射.

对每个 $\alpha\in H(b,c)$,
\[
P_{\alpha}(X)\overset{\text{def}}{=}(X-\alpha)(X-\bar{\alpha})=X^{2}-(\alpha+\bar{\alpha})X+\Nm(\alpha)\in K[X].
\]
因此 $P_{\alpha}(X)$ 是 $\alpha$ 在扩张 $K[\alpha]/K$ 中的特征多项式. 特别地,
\[
\Nm(\alpha)=\Nm_{K[\alpha]/K}(\alpha).
\]
现在 (见 (IV, \ref{rem:iv.4.4})) $H(b,c)$ 包含 $K$ 的每个二次扩张的同构副本, 例如 $K$ 的非分歧二次扩张与一个全分歧二次扩张. 因此 $\Nm(H(b,c)^{\times})$ 包含这两个互异二次扩张的范数群, 从而 (如上) 等于 $K^{\times}$.
\end{proof}

\begin{corollary}\label{cor:viii.3.10}
$\Q_p$ 的有限扩张上 $\ge 5$ 元的每个非退化二次型都表示 $0$.
\end{corollary}

\begin{proof}
可写 $q=r(X_1,\ldots,X_4)-aX_5^{2}+q'(X_6,\ldots)$, 其中 $r$ 是四元非退化二次型, $a\ne 0$. 则 $r$ 表示 $a$, 故 $q$ 表示 $0$.
\end{proof}

\begin{corollary}\label{cor:viii.3.11}
数域 $K$ 上 $\ge 5$ 元的非退化二次型 $q$ 表示 $0$ 当且仅当它在 $K$ 的每个实完备化中表示 $0$.
\end{corollary}

\begin{proof}
结合 (\ref{cor:viii.3.10}) 与 (\ref{thm:viii.3.5}).
\end{proof}

\begin{remark}\label{rem:viii.3.12}
命题\ref{prop:viii.3.9} 的证明对 $K=\R$ 也成立, 直到最后一步: $\R$ 的唯一的二次扩张是 $\CC$, 故
\[
\Nm(\mathbb{H}(b,c)^{\times})=\Nm(\CC^{\times})=\R_{>0}.
\]
因此它表明: $\R$ 上不表示零的四元非退化型表示所有严格正的实数.
\end{remark}
\section{基本正合列与基本类}\label{sec:VIII.4}

对 Galois 扩张 $L/K$, 记
\[
\mathrm{Br}(L/K)=H^{2}(\Gal(L/K),L^{\times});
\]
对数域的 Galois 扩张 $L/K$, 记
\[
H^{2}(L/K)=H^{2}(\Gal(L/K),\mathbf{C}_{L}).
\]
由于 $H^{1}(G,L^{\times})=0$ (见 (II, \ref{prop:ii.1.22})) 且 $H^{1}(G,\mathbf{C}_{L})=0$ (见 (VII, \ref{thm:vii.5.1})), 对每个 Galois 扩张塔 $E\supset L\supset K$, 我们得到正合列
\[
0\to\mathrm{Br}(L/K)\to\mathrm{Br}(E/K)\to\mathrm{Br}(E/L)
\]
与
\[
0\to H^{2}(L/K)\to H^{2}(E/K)\to H^{2}(E/L).
\]
对越来越大的域 $E\subset K^{\mathrm{al}}$ 取正向极限, 得到正合列
\[
0\to\mathrm{Br}(L/K)\to\mathrm{Br}(K)\to\mathrm{Br}(L)
\]
与
\[
0\to H^{2}(L/K)\to H^{2}(K^{\mathrm{al}}/K)\to H^{2}(K^{\mathrm{al}}/L).
\]
于是, 我们可以把 $\mathrm{Br}(L/K)$ 视为 $\mathrm{Br}(K)$ 中被 $L$ 分裂的元素构成的子群; 对 $H^{2}(L/K)$ 亦然.

设 $L/K$ 是 $n$ 次有限 Galois 的数域扩张, 并考虑图表:
\begin{equation*}
\begin{tikzcd}[column sep=1.2em, row sep=large]
& & & H^{2}(L/K) \\
0 \ar[r] & \mathrm{Br}(L/K) \ar[r] & \bigoplus_{v}\mathrm{Br}(L^{v}/K_{v}) \ar[ur] \ar[dr, "\varSigma"'] & \\
& & & \Q/\Z
\end{tikzcd}
\end{equation*}
上 行是
\[
0\to L^{\times}\to\mathbb{I}_{L}\to\mathbf{C}_{L}\to0
\]
的上同调列的一部分. 左端的零来自 $H^{1}(G,\mathbf{C}_L)=0$. 上 行正合, 但映射 $\bigoplus_v\mathrm{Br}(L^{v}/K_v)\to H^{2}(L/K)$ 一般不是满射——我们把它记为 $H^{2}(L/K)'$.

回顾 (III, \ref{thm:iii.2.1}): 对每个素 $v$, 有同态
\[
\mathrm{inv}_{v}:\mathrm{Br}(K_{v})\to\Q/\Z.
\]
若 $v$ 非阿基米德, 它是 $\mathrm{Br}(K_v)$ 到 $\Q/\Z$ 的同构; 若 $v$ 是实的, 它是 $\mathrm{Br}(K_v)$ 到 $\frac{1}{2}\Z/\Z$ 的同构. 此外, 若 $L_w/K_v$ 的次数为 $n_v$, 则 $\mathrm{inv}_w(\beta)=n_v\cdot\mathrm{inv}_v(\beta)$, 故 $\mathrm{inv}_v$ 定义同构
\[
\mathrm{inv}_{v}:\mathrm{Br}(L_{w}/K_{v})\to\tfrac{1}{n_{v}}\Z/\Z.
\]
图表中的东南向箭头是
\[
\varSigma:\bigoplus_{v}\mathrm{Br}(L^{v}/K_{v})\to\Q/\Z,\qquad(\beta_{v})\mapsto\sum\mathrm{inv}_{v}(\beta_{v}).
\]
$\mathrm{inv}_v$ 的像是循环群 $\frac{1}{n}\Z/\Z$ 中 $n_v$ 阶的循环子群, 故 $\varSigma$ 的像是循环子群 $\frac{1}{n_0}\Z/\Z$, 其中 $n_0=\mathrm{lcm}(n_v)$.

根据定理 (VII, \ref{thm:vii.5.1}), $H^{2}(L/K)$ 的阶整除 $n$. 根据定理 (VII, \ref{thm:vii.8.1}), 图表的下 行是复形, 故图表中的映射诱导满同态
\[
H^{2}(L/K)^{\prime}\to\tfrac{1}{n_{0}}\Z/\Z.
\]
局部次数的最小公倍数 $n_0$ 总是整除 $n$, 但不必等于 $n$ (见下面的例\ref{ex:viii.4.5}). 然而, 假设扩张 $L/K$ 具有性质 $n=n_0$. 则:

(a) 映射 $H^{2}(L/K)'\to\frac{1}{n_0}\Z/\Z$ 是同构. (它是满射, 且 $(H^{2}(L/K)':1)\le(H^{2}(L/K):1)\le n$.)

(b) $H^{2}(L/K)'=H^{2}(L/K)$, 且两者的阶都是 $n$.

(c) 下 行是正合列
\[
0\to\mathrm{Br}(L/K)\to\bigoplus_{v}\mathrm{Br}(L^{v}/K_{v})\xrightarrow{\ \varSigma\ }\tfrac{1}{n}\Z/\Z\to0
\]
(因为它同构于上 行).

\begin{lemma}\label{lem:viii.4.1}
若 $L/K$ 是循环的, 则 $n=n_0$.
\end{lemma}

\begin{proof}
设 $S\supset S_\infty$ 是包含所有在 $L$ 中分歧的素的 $K$ 的素集. 对 $v\notin S$, $(p_v,L/K)$ 是 $\Gal(L/K)$ 中 $n_v$ 阶 ($=f_v$) 的元素, 故 Artin 映射 $I^{S}\to\Gal(L/K)$ 的像的阶为 $n_0\overset{\text{def}}{=}\mathrm{lcm}(n_v)$. 根据互反律 (VII, \ref{cor:vii.4.8}), Artin 映射是满射, 这蕴含 $n_0=n$. [利用复分析可以证明更多: 对密度为 $\varphi(n)/n$ 的一组 $v$, $L^{v}/K_v$ 是 $n$ 阶循环的: 设 $\mathfrak{m}$ 为 $L/K$ 的模, 并取 $\mathfrak{a}\in I^{S}$ 使 $(\mathfrak{a},L/K)$ 生成 $\Gal(L/K)$; 则在 $C_m$ 中同余于 $\mathfrak{a}$ 的素理想集的密度为 $1/n$.]
\end{proof}

设 $\Q^{\mathrm{c}}$ 是 (I, 例 \ref{ex:A.5}(d)) 中定义的 $\Q$ 的无限循环分圆扩张 (另见 (VII, \ref{lem:vii.7.3})), 并令 $\Omega=\Q^{\mathrm{c}}\cdot K$. 对每个 $n$, $\Omega$ 包含 $K$ 的唯一的 $n$ 次循环扩张 $\Omega_n$. 前面的引理与说明表明
\[
0\to\mathrm{Br}(\Omega_{n}/K)\to\bigoplus_{v}\mathrm{Br}(\Omega_{n}^{v}/K_{v})\xrightarrow{\ \varSigma\ }\tfrac{1}{n}\Z/\Z\to0
\]
正合. 对所有 $n$ 取正向极限 (实际上是定向并), 得到正合列
\[
0\to\mathrm{Br}(\Omega/K)\to\bigoplus_{v}\mathrm{Br}(\Omega^{v}/K_{v})\to\Q/\Z\to0.
\]

\begin{theorem}\label{thm:viii.4.2}
对每个数域 $K$, 序列
\[
0\to\mathrm{Br}(K)\to\bigoplus_{v}\mathrm{Br}(K_{v})\xrightarrow{\ \varSigma\ }\Q/\Z\to0
\]
正合.
\end{theorem}

\begin{proof}
根据命题 (VII, \ref{prop:vii.7.2}), $\mathrm{Br}(\Omega/K)=\mathrm{Br}(K)$. 而且 $\mathrm{Br}(\Omega^{v}/K_v)=\mathrm{Br}(K_v)$, 因为在非阿基米德情形, 当 $n\to\infty$ 时 $[\Omega_n^{v}:K_v]\to\infty$ (见 (VII, \ref{lem:vii.7.3})).
\end{proof}

定理中的序列称为整体类域论的\textbf{基本正合列}\index{jibenzhenghelie@基本正合列(fundamental exact sequence)}.

\begin{corollary}\label{cor:viii.4.3}
对每个有限扩张 $L/K$, 序列
\[
0\to\mathrm{Br}(L/K)\to\bigoplus_{v}\mathrm{Br}(L^{v}/K_{v})\xrightarrow{\ \varSigma\ }\tfrac{1}{n_{0}}\Z/\Z\to0,\qquad n_{0}=\mathrm{lcm}(n_{v}),
\]
正合.
\end{corollary}

\begin{proof}
把 $K$ 的基本正合列映到 $L$ 的基本正合列, 对所得图表应用蛇形引理.
\end{proof}

\begin{example}\label{ex:viii.4.4}
(a) 对数域的有限循环扩张 $L/K$, 基本正合列可等同于
\[
0\to K^{\times}/\Nm(L^{\times})\to\bigoplus_{v}K_{v}^{\times}/\Nm(L^{v\times})\to\tfrac{1}{n}\Z/\Z\to0,\qquad n=[L:K].
\]

(b) 设 $D$ 是数域 $K$ 上的可除代数, 并设 $i_v=\mathrm{inv}_v(D\otimes_K K_v)$. 则: $i_v=0$ 对除有限多个之外的所有 $v$; $i_v=0$ 若 $v$ 为复; $i_v\in\frac{1}{2}\Z/\Z$ 若 $v$ 为实; 且 $\sum i_v\equiv 0 \bmod \Z$. 族 $(i_v)$ 确定 $D$ 的同构类, 而任何满足这些条件的族 $(i_v)$ 都是某个可除代数的不变量族. 显然, $D$ 的类在 $\mathrm{Br}(K)$ 中的阶是诸 $i_v$ 的最小公分母 $n$. 还可以证明 $[D:K]=n^{2}$. 例如, 给出 $\Q$ 上的一个四元数代数就等价于给出 $\Q$ 的一个由偶数个元素构成的有限素集.
\end{example}

\begin{example}\label{ex:viii.4.5}
设 $L=\Q[\sqrt{13},\sqrt{17}]$. 显然 $n=4$, 但我断言对所有 $v$ 有 $n_v=1$ 或 $2$. 由于 $13$ 与 $17$ 都同余于 $1$ 模 $4$, $2$ 在 $L$ 中不分歧. 因此, 对 $w\mid p$, $p\ne 13,17$, $L_w$ 是 $\Q_p$ 的非分歧扩张. 特别地, 其 Galois 群是循环的. 由于它是 $\Gal(L/\Q)$ 的子群, 它被 $2$ 杀死, 从而阶为 $1$ 或 $2$. 另一方面, $\left(\frac{17}{13}\right)=1$ (显然) 且 $\left(\frac{13}{17}\right)=\left(\frac{17}{13}\right)=1$. 故 $17$ 模 $13$ 是平方, 而 Hensel 引理表明它在 $\Q_{13}$ 中是平方. 同样, $13$ 在 $\Q_{17}$ 中是平方.
\end{example}

\subsection*{基本类}

由以上讨论可知, 存在同构
\[
\mathrm{inv}_{K}:H^{2}(\Omega/K)\to\Q/\Z
\]
由如下性质唯一刻画: 复合
\[
\bigoplus_{v}\mathrm{Br}(K_{v})\to H^{2}(\varOmega/K)\xrightarrow{\ \mathrm{inv}_{K}\ }\Q/\Z
\]
是 $(\beta_{v})\mapsto\sum\mathrm{inv}_{v}(\beta_{v})$.

\begin{lemma}\label{lem:viii.4.6}
对数域的每个有限扩张 $L/K$ 与每个 $\gamma\in H^{2}(\Omega/K)$, 有 $\mathrm{inv}_L(\gamma)=n\,\mathrm{inv}_K(\gamma)$, 其中 $n=[L:K]$.
\end{lemma}

\begin{proof}
利用局部次数之和等于整体次数这一事实.
\end{proof}

因此, 对每个有限 Galois 的 $L/K$, 得到同构
\[
\mathrm{inv}_{L/K}:H^{2}(L/K)\to\tfrac{1}{n}\Z/\Z.
\]
对所有 $L\subset K^{\mathrm{al}}$ 取正向极限, 得到同构
\[
\mathrm{inv}_{K}:H^{2}(K^{\mathrm{al}}/K)\to\Q/\Z.
\]

\begin{theorem}\label{thm:viii.4.7}
对数域的每个有限 Galois 扩张 $L/K$, $H^{2}(L/K)$ 是 $n=[L:K]$ 阶的循环群, 带有典范生成元 $u_{L/K}$.
\end{theorem}

\begin{proof}
取 $u_{L/K}$ 为满足 $\mathrm{inv}_{L/K}(u_{L/K})=\frac{1}{n}\bmod\Z$ 的元素.
\end{proof}

$H^{2}(L/K)$ 的生成元 $u_{L/K}$ 称为\textbf{基本类}\index{jibenlei@基本类(fundamental class)}. 可以像局部情形 (见 (III, \ref{lem:iii.2.7})) 一样证明: 对每个由有限 Galois 扩张组成的塔 $E\supset L\supset K$,
\[
\begin{aligned}
\Res(u_{E/K})&=u_{E/L},\\
\Inf(u_{L/K})&=[E:L]\,u_{E/K}.
\end{aligned}
\]
因此, 可以应用 Tate 定理 ((II, \ref{thm:ii.3.11})) 得到同构
\[
\Gal(L/K)^{\mathrm{ab}}\to\mathbf{C}_{K}/\Nm(\mathbf{C}_{L}).
\]
该同构与上一章定义的整体 Artin 映射 $\phi_{L/K}$ 互逆, 这一点由整体基本类与局部基本类相容这一事实推出.

\subsection*{范数限制定理}

\begin{theorem}\label{thm:viii.4.8}
设 $E$ 是 $K$ 的有限扩张 (不必 Galois), 并设 $M$ 是 $E$ 的使 $M/K$ 为 Abel Galois 扩张的最大子扩张. 则
\[
\Nm_{E/K}\mathbf{C}_{E}=\Nm_{M/K}\mathbf{C}_{M}.
\]
\end{theorem}

\begin{proof}
设 $L$ 是 $K$ 的包含 $E$ 的 Galois 扩张, 并令 $G=\Gal(L/K)$, $H=\Gal(L/E)$. 考虑交换图表
\[
\begin{tikzcd}[column sep=large]
H_{T}^{-2}(H,\Z) \ar[r, "\simeq"] \ar[d, "\mathrm{Cor}"'] & H_{T}^{0}(H,\mathbf{C}_{L}) \ar[d, "\mathrm{Cor}"] \\
H_{T}^{-2}(G,\Z) \ar[r, "\simeq"'] & H_{T}^{0}(G,\mathbf{C}_{L})
\end{tikzcd}
\]
其中的水平箭头是与基本类的杯积. 它可以等同于交换图表:
\[
\begin{tikzcd}[column sep=large]
H^{\mathrm{ab}} \ar[r, "\simeq"] \ar[d] & C_{E}/\Nm_{L/E}C_{L} \ar[d] \\
G^{\mathrm{ab}} \ar[r, "\simeq"'] & C_{K}/\Nm_{L/K}C_{L}.
\end{tikzcd}
\]
故 $H^{\mathrm{ab}}\to G^{\mathrm{ab}}$ 的余核同构于 $C_K/\Nm_{E/K}(C_E)$. 但该余核等于 $\Gal(M/K)$, 后者又同构于 $C_K/\Nm_{M/K}(C_M)$. 由于 $\Nm(C_M)\supset\Nm(C_E)$, 这两个群必相等.
\end{proof}
\section{高次互反律}\label{sec:VIII.5}

回顾: 对奇素数 $p$ 与不被 $p$ 整除的整数 $a$, Legendre 符号 (或二次剩余符号)
\[
\left(\frac{a}{p}\right)\overset{\text{def}}{=}
\begin{cases}
1 & \text{若 } a \text{ 模 } p \text{ 是平方},\\
-1 & \text{否则}
\end{cases}
\]
的定义. 群 $\F_p^{\times}$ 是 $p-1$ 阶循环群, 以 $-1$ 为其唯一的 $2$ 阶元素. 因此, 对 $u\in\F_p^{\times}$, $u^{\frac{p-1}{2}}$ 依 $u$ 是否为平方而等于 $1$ 或 $-1$, 故 $\left(\frac{a}{p}\right)$ 是满足
\[
\left(\frac{a}{p}\right)\equiv a^{\frac{p-1}{2}}\pmod p
\]
的 $1$ 的唯一平方根.

二次互反律说: 对奇素数 $p,q$,
\[
\left(\frac{p}{q}\right)\left(\frac{q}{p}\right)=(-1)^{\frac{p-1}{2}\cdot\frac{q-1}{2}}.
\]
二次互反律的补充说
\[
\left(\frac{-1}{p}\right)=(-1)^{\frac{p-1}{2}},\qquad
\left(\frac{2}{p}\right)=(-1)^{\frac{n^{2}-1}{8}}.
\]

对 Gauss 整数 $\alpha$ (即 $\Z[i]$ 的元素) 与奇 Gauss 素数 $\pi$ (即 $\Z[i]$ 中不整除 $2$ 的素元), Gauss 定义 $\left(\frac{\alpha}{\pi}\right)$ (四次剩余符号) 为满足
\[
\left(\frac{\alpha}{\pi}\right)\equiv\alpha^{\frac{\mathbb{N}\pi-1}{4}}\pmod{\pi}
\]
的唯一的 $1$ 的 $4$ 次根, 并对这些符号证明了一个四次互反律. 后来 Eisenstein 证明了一个三次互反律. Emil Artin 指出, 他的定理 (V, \ref{thm:v.3.5}) 蕴含所有可能的此类互反律, 因而可以视为``不含 $n$ 次单位根的域的互反律''. 本节其余部分解释这一评论.

\subsection*{幂剩余符号}

设 $K$ 是包含本原 $n$ 次单位根的数域. 对 $K$ 的元素的每个有限组 $a,b,\ldots$, 定义 $S(a,b,\ldots)$ 为满足下列条件之一的 $K$ 的素理想集: $\mathrm{ord}_{\mathfrak{p}}(n)\ne 0$, 或 $\mathrm{ord}_{\mathfrak{p}}(a)\ne 0$, 或 $\mathrm{ord}_{\mathfrak{p}}(b)\ne 0$, \ldots. 特别地, $S$ 本身仅由 $n$ 的因子组成.

回顾: $X^{n}-1$ 的判别式只被整除 $n$ 的素整除. 因此, 对任何 $\mathfrak{p}\nmid n$, $X^{n}-1$ 在 $\F_p^{\mathrm{al}}$ 中有 $n$ 个互异的根, 且映射
\[
\zeta\mapsto\zeta \bmod \mathfrak{p}:\mu_{n}(K)\to\mu_{n}(\mathcal{O}_{K}/\mathfrak{p})
\]
对任何满足 $\mathfrak{p}\nmid n$ 的素理想是双射. 对这样的素 $\mathfrak{p}$, 令 $q=\mathbb{N}\mathfrak{p}\overset{\text{def}}{=}(\mathcal{O}_{K}:\mathfrak{p})$. 则 $\F_q^{\times}$ 是 $q-1$ 阶循环群, 故 $n|q-1$, 且 $\zeta^{\frac{q-1}{n}}\in\mu_n\subset\F_q^{\times}$.

对 $a\in K^{\times}$ 与 $\mathfrak{p}\notin S(a)$, 定义 $\left(\frac{a}{\mathfrak{p}}\right)$ 为满足
\[
\left(\frac{a}{\mathfrak{p}}\right)\equiv a^{\frac{\mathbb{N}\mathfrak{p}-1}{n}}\pmod{\mathfrak{p}}
\]
的唯一的 $1$ 的 $n$ 次根.

\noindent\textbf{5.1}\quad 对所有 $a,b\in K^{\times}$ 与 $\mathfrak{p}\notin S(a,b)$,
\[
\left(\frac{ab}{\mathfrak{p}}\right)=\left(\frac{a}{\mathfrak{p}}\right)\left(\frac{b}{\mathfrak{p}}\right).
\]
由定义显然.

\noindent\textbf{5.2}\quad 对 $a\in K^{\times}$ 与 $\mathfrak{p}\notin S(a)$, 下述等价:

(a) $\left(\frac{a}{\mathfrak{p}}\right)=1$;

(b) $a$ 在 $\mathcal{O}_{K}/\mathfrak{p}$ 中成为 $n$ 次幂;

(c) $a$ 在 $K_{\mathfrak{p}}$ 中成为 $n$ 次幂.

(a) 与 (b) 的等价性由
\[
1\to\F_{q}^{\times n}\to\F_{q}^{\times}\xrightarrow{\ x\mapsto x^{\frac{q-1}{n}}\ }\mu_{n}\to1,\qquad q=\mathbb{N}\mathfrak{p}
\]
的正合性推出. 若 $X^{n}-a$ 模 $\mathfrak{p}$ 有解, 则 Hensel 引理 (ANT, 7.33) 表明它在 $K_{\mathfrak{p}}$ 中有解. 反之, 若 $a=\alpha^{n}$, $\alpha\in K_{\mathfrak{p}}$, 则 $\mathrm{ord}_{\mathfrak{p}}(\alpha)=\frac{1}{n}\mathrm{ord}_{\mathfrak{p}}(a)=0$, 故 $\alpha\in\mathcal{O}_{K_{\mathfrak{p}}}$. 映射 $\mathcal{O}_K\to\mathcal{O}_{K_{\mathfrak{p}}}/\mathfrak{p}$ 是满射, 故存在 $\alpha_0\in\mathcal{O}_K$ 映到 $\alpha$ 模 $\mathfrak{p}$ 的像.

我们用线性性把映射 $\mathfrak{p}\mapsto\left(\frac{a}{\mathfrak{p}}\right)$ 延拓到 $I^{S(a)}$: 即对 $\mathfrak{b}=\prod\mathfrak{p}_i^{r_i}\in I^{S(a)}$,
\[
\left(\frac{a}{\mathfrak{b}}\right)=\prod\left(\frac{a}{\mathfrak{p}_{i}}\right)^{r_{i}}.
\]
我们把 $\left(\frac{a}{\mathfrak{b}}\right)$ 简写为 $\left(\frac{a}{b}\right)$.

对使 $S'$ 中的素都不分歧的 Abel 扩张 $L/K$,以 $\psi_{L/K}\colon I^{S}\to\Gal(L/K)$ 表示 Artin 映射 (见第 \ref{ch:V} 章).

\noindent\phantomsection\hypertarget{item:viii.5.3}{}\textbf{5.3}\quad 对所有 $a\in K^{\times}$ 与 $\mathfrak{b}\in I^{S(a)}$,
\[
\psi_{K[a^{\frac{1}{n}}]/K}(\mathfrak{b})\big(a^{\frac{1}{n}}\big)=\left(\frac{a}{\mathfrak{b}}\right)a^{\frac{1}{n}}.
\]
由 Galois 理论, 存在 $1$ 的 $n$ 次根 $\zeta(\mathfrak{b})$ 使 $\psi(\mathfrak{b})(a^{\frac{1}{n}})=\zeta(\mathfrak{b})\cdot a^{\frac{1}{n}}$, 且映射 $\mathfrak{b}\mapsto\zeta(\mathfrak{b})$ 是同态. 因此只需对 $\mathfrak{b}=\mathfrak{p}$ 为素理想的情形证明等式. 按定义,
\[
\psi(\mathfrak{p})(x)\equiv x^{\mathbb{N}\mathfrak{p}}\pmod{\mathfrak{p}}.
\]
由
\[
\psi(\mathfrak{p})\big(a^{\frac{1}{n}}\big)=\zeta(\mathfrak{p})\cdot a^{\frac{1}{n}}
\]
得到
\[
\zeta(\mathfrak{p})\cdot a^{\frac{1}{n}}\equiv a^{\frac{\mathbb{N}\mathfrak{p}}{n}}\pmod{\mathfrak{p}},
\]
由此推出 $\zeta(\mathfrak{p})=\left(\frac{a}{\mathfrak{p}}\right)$.

\noindent\textbf{5.4}\quad 设 $a\in\mathcal{O}_{K}$, 并设 $\mathfrak{b}$ 是 $I^{S(a)}$ 中的整理想. 若 $a'\in\mathcal{O}_{K}$ 满足 $a'\equiv a \bmod \mathfrak{b}$, 则 $\mathfrak{b}\in I^{S(a')}$ 且
\[
\left(\frac{a}{\mathfrak{b}}\right)=\left(\frac{a'}{\mathfrak{b}}\right).
\]
事实上, 对每个整除 $\mathfrak{b}$ 的素理想 $\mathfrak{p}$, 有 $a'\equiv a \bmod \mathfrak{p}$, 故 $\left(\frac{a}{\mathfrak{p}}\right)=\left(\frac{a'}{\mathfrak{p}}\right)$.

Artin 互反律使我们能为 $\left(\frac{a}{\mathfrak{b}}\right)$ 作为 $\mathfrak{b}$ 的函数证明一个类似但较弱的结果.

\noindent\phantomsection\hypertarget{item:viii.5.5}{}\textbf{5.5}\quad 设 $a\in K^{\times}$. 则存在支撑在 $S(a)$ 中的模 $\mathfrak{m}$, 使得 $\left(\frac{a}{\mathfrak{b}}\right)$ 只依赖于 $\mathfrak{b}$ 在射线类群 $C_\mathfrak{m}$ 中的类.

根据命题 (VII, \ref{prop:vii.A.5}), $S(a)$ 包含所有在 $K[a^{\frac{1}{n}}]$ 中分歧的素. 因此 Artin 互反律 (V, \ref{thm:v.3.5}) 表明存在支撑在 $S(a)$ 中的模 $\mathfrak{m}$, 使得 $\psi(\mathfrak{b})$ 只依赖于 $\mathfrak{b}$ 在射线类群 $C_\mathfrak{m}$ 中的类.

\begin{NOTES}
抄自 \mo{299588}. 幂剩余符号是 $K$ 的一个酉 Hecke 特征, 无穷型平凡, 阶有限, 导子整除 $\mathfrak{m}$ 的某个幂. 这一事实可由类域论推出, 例如由本讲义 \hyperlink{item:viii.5.5}{VIII, 5.5}. (解释待补.)
\end{NOTES}

\subsection*{Hilbert 符号}

设 $K_v$ 是包含本原 $n$ 次单位根的局部域. Hilbert 符号是一个配对
\[
a,b\mapsto(a,b)_{v}:K_{v}^{\times}/K_{v}^{\times n}\times K_{v}^{\times}/K_{v}^{\times n}\to\mu_{n},
\]
其中 $\mu_{n}$ 是 $K_{v}$ 中 $1$ 的 $n$ 次根构成的群. 定义它的最自然方式也许是杯积映射
\[
H^{1}(G,\mu_{n})\times H^{1}(G,\mu_{n})\to H^{2}(G,\mu_{n}\otimes\mu_{n}),\qquad G=\Gal(K^{\mathrm{al}}/K),
\]
再接上由不变映射 $\mathrm{inv}_v$ 定义的同构
\[
H^{2}(G,\mu_{n}\otimes\mu_{n})=H^{2}(G,\mu_{n})\otimes\mu_{n}\to\mu_{n}.
\]
不过, 出于 1920 与 1930 年代的精神, 我将用中心单代数来定义它.

回顾 (IV, \S \ref{sec:IV.5}): 对任意 $a,b\in K_v^{\times}$, 定义 $A(a,b;\zeta)$ 为 $K_v$-代数, 它由元素 $i,j$ 生成, 满足关系
\[
i^{n}=a,\qquad j^{n}=b,\qquad ij=\zeta ji.
\]
它是 $K_v$ 上次数为 $n$ 的中心单代数. 当 $n=2$ 时, $A(a,b;-1)$ 就是四元数代数 $H(a,b)$. 我们定义
\[
(a,b)_{v}=\zeta^{-n\cdot\mathrm{inv}_{v}([A(a,b;\zeta)])},
\]
其中 $[A(a,b;\zeta)]$ 是 $A(a,b;\zeta)$ 在 $\mathrm{Br}(K_v)$ 中的类. 由于 $A(a,b;\zeta)$ 被一个 $n$ 次域分裂 (例如, 在 $n=2$ 的情形, $\Q[i]$ 分裂 $A(a,b;-1)$), 其不变量属于 $\frac{1}{n}\Z/\Z$, 故 $n\cdot\mathrm{inv}_v([A(a,b;\zeta)])$ 是 $\Z/n\Z$ 的元素. 显然 $A(a,b;\zeta)$ 的同构类只依赖于作为 $K_v^{\times}/K_v^{\times n}$ 中元素的 $a$ 与 $b$, 因此我们确实得到配对
\[
K_{v}^{\times}/K_{v}^{\times n}\times K_{v}^{\times}/K_{v}^{\times n}\to\mu_{n}.
\]
不过, 从这个视角并不能明显看出该配对是双线性的.

\setcounter{theorem}{5}
\begin{example}\label{ex:viii.5.6}
考虑 $K_v=\Q_p$ ($p$ 为奇素数), $n=2$ 的情形. 则 $(a,b)_p=\pm 1$, 且
\[
\begin{aligned}
(a,b)_{p}=1&\iff H(a,b)\approx M_{2}(K_{v});\\
&\iff X^{2}-aY^{2}-bZ^{2}+abT^{2} \text{ 在 } K_{v} \text{ 中表示 } 0;\\
&\iff b \text{ 是来自 } K[\sqrt{a}] \text{ 的范数};\\
&\iff X^{2}-aY^{2}-bZ^{2} \text{ 在 } K_{v} \text{ 中表示 } 0.
\end{aligned}
\]
证明这些等价时分别用到: 四元数代数的不变量为 $\frac{1}{2}$ 当且仅当它是可除代数; 习题 IV.\ref{ex:iv.5.7}; 命题\ref{prop:viii.3.7}(d); 命题\ref{prop:viii.3.7}(c). 最后一个条件表明我们对 Hilbert 符号的定义与例如 \cite{Serre1970}, III 中的定义一致.
\end{example}

\noindent\textbf{5.7}\quad 对任意 $a,b$,
\[
A(b,a;\zeta)\approx A(a,b;\zeta^{-1})\approx A(a,b;\zeta)^{\mathrm{opp}}.
\]
因此
\[
(b,a)_{v}=(a,b)_{v}^{-1}.
\]
按定义, $A(b,a;\zeta)$ 是由 $i',j'$ 生成的 $K_v$-代数, 满足关系 $i'^{n}=b$, $j'^{n}=a$, $i'j'=\zeta j'i'$. 映射 $i'\mapsto j$, $j'\mapsto i$ 是同构 $A(b,a;\zeta)\to A(a,b;\zeta^{-1})$. 映射 $i\mapsto i$, $j\mapsto j$ 是同构 $A(a,b;\zeta)^{\mathrm{opp}}\to A(a,b;\zeta^{-1})$.

\noindent\phantomsection\hypertarget{item:viii.5.8}{}\textbf{5.8}\quad 设 $a,b\in K^{\times}$. 对任意 $v\notin S(a)$,
\[
(a,b)_{v}=\left(\frac{a}{\mathfrak{p}_{v}}\right)^{\mathrm{ord}_{v}(b)}.
\]
为简单起见, 设 $A(a,b;\zeta)$ 是可除代数. 回顾 (IV, \S \ref{sec:IV.4}): 为计算局部域 $K_v$ 上中心可除代数 $D$ 的不变量, 我们

(a) 选取极大非分歧子域 $L\subset D$;

(b) 找出元素 $\beta\in D$ 使 $\alpha\mapsto\beta\alpha\beta^{-1}$ 是 $L$ 的 Frobenius 自同构 (这样的 $\beta$ 由 Noether--Skolem 定理保证存在);

(c) 令 $\mathrm{inv}_{v}([D])=\mathrm{ord}_{v}(\beta)$.

我们把它应用于 $L=K_v[i]=K_v[a^{\frac{1}{n}}]$. 注意, 由于 $v\notin S(a)$, 该扩张非分歧. 设 $\left(\frac{a}{\mathfrak{p}_v}\right)=\zeta^{r}$, 即 $(\mathfrak{p},L/K_v)(i)=\zeta^{r}i$. 由于 $jij^{-1}=\zeta^{-1}i$, 可取 $\beta=j^{-r}$. 则 $\beta^{n}=b^{-r}$, 故 $\mathrm{ord}_v(\beta)=-\frac{r}{n}\mathrm{ord}_v(b)$. 于是
\[
(a,b)_{v}\overset{\text{def}}{=}\zeta^{-n\,\mathrm{inv}_{v}(A(a,b;\zeta))}=\zeta^{r\cdot\mathrm{ord}_{v}(b)}=\left(\frac{a}{\mathfrak{p}_{v}}\right)^{\mathrm{ord}_{v}(b)}.
\]

\setcounter{theorem}{8}
\begin{remark}\label{rem:viii.5.9}
事实上,
\[
(a,b)_{v}=\frac{\phi_{v}(b)\big(a^{\frac{1}{n}}\big)}{a^{\frac{1}{n}}}
\]
对所有 $a,b,v$ 成立. 见 (III, \ref{rem:iii.4.5}).
\end{remark}

\noindent\textbf{5.10}\quad 对 $a,b\in K^{\times}$,
\[
\prod_{v}(a,b)_{v}=1.
\]
在证明互反律的过程中, 我们证明了对每个 $\beta\in\mathrm{Br}(K)$, $\sum\mathrm{inv}_v(\beta)=0$. 特别地, $\sum\mathrm{inv}_v(A(a,b;\zeta))=0$, 这就蕴含了该公式.

对 $a,b\in K^{\times}$, 定义
\[
\left(\frac{a}{b}\right)=\prod_{v\notin S(a)}\left(\frac{a}{v}\right)^{\mathrm{ord}_{v}(b)}=\left(\frac{a}{(b)^{S(a)}}\right),
\]
其中 $(b)^{S(a)}$ 是由 $b$ 生成的 $I^{S(a)}$ 中的理想. 符号 $\left(\frac{a}{b}\right)$ 对 $b$ 是乘性的, 但 $\left(\frac{aa'}{b}\right)=\left(\frac{a}{b}\right)\left(\frac{a'}{b}\right)$ 一般不成立, 除非 $S(b)\cap S(a,a')=S$.

\setcounter{theorem}{10}
\begin{theorem}[幂互反律]\label{thm:viii.5.11}
设 $a,b\in K^{\times}$ 满足 $S(a)\cap S(b)=S$ (例如 $a$ 与 $b$ 可以互素). 则
\[
\left(\frac{a}{b}\right)\left(\frac{b}{a}\right)^{-1}=\prod_{v\in S}(b,a)_{v}.
\]
此外, 若 $S(c)=S$, 则
\[
\left(\frac{c}{b}\right)=\prod_{v\in S}(b,c)_{v}.
\]
\end{theorem}

\begin{proof}
令 $S'(a)=S(a)\setminus S$, $S'(b)=S(b)\setminus S$. 我们的假设是 $S'(a)$ 与 $S'(b)$ 不相交. 则
\[
\left(\frac{a}{b}\right)=\prod_{v\in S'(b)}\left(\frac{a}{\mathfrak{p}_{v}}\right)^{\mathrm{ord}_{v}(b)}=\prod_{v\in S'(b)}(a,b)_{v}
\]
且
\[
\left(\frac{b}{a}\right)=\prod_{v\in S'(a)}\left(\frac{b}{\mathfrak{p}_{v}}\right)^{\mathrm{ord}_{v}(a)}=\prod_{v\in S'(a)}(b,a)_{v}.
\]
因此
\[
\left(\frac{a}{b}\right)\left(\frac{b}{a}\right)^{-1}=\prod_{v\in S'(a)\cup S'(b)}(a,b)_{v}.
\]
对 $v\notin S\cup S'(a)\cup S'(b)$, $(a,b)_v=1$ (例如由 \hyperlink{item:viii.5.8}{5.8}), 故乘积公式表明
\[
\prod_{v\in S'(a)\cup S'(b)}(a,b)_{v}\times\prod_{v\in S}(a,b)_{v}=1.
\]
这就完成了第一个等式的证明; 第二个是显然的.
\end{proof}

要得到完全显式的公式, 还需计算 $v\in S$ 时的 Hilbert 符号. 对无限素这很容易: 若 $v$ 为复, 则 $(a,b)_v=1$ 恒成立; 若 $v$ 为实, 则
\[
(a,b)_{v}=1\iff X^{2}-aY^{2}-bZ^{2} \text{ 表示 } 0\iff a>0 \text{ 或 } b>0.
\]
对 $K=\Q$ 与 $n=2$,
\[
(u2^{r},v2^{s})_{2}=(-1)^{\frac{u-1}{2}\cdot\frac{v-1}{2}+r\frac{v^{2}-1}{8}+s\frac{u^{2}-1}{8}},
\]
其中 $u,v$ 是 2-adic 单位, 指数模 $2$ 理解. 其初等证明见 \cite{Serre1970}, III, 1.2. 依次把这个公式应用于 $(p,q)$ ($p,q$ 为奇素数)、$(2,p)$ ($p$ 为奇素数) 与 $(-1,p)$ ($p$ 为奇素数), 便得到经典的二次互反律 (包括其补充).

对奇素数 $p$ 与 $K=\Q[\zeta]$ ($\zeta$ 为本原 $p$ 次单位根), 可以把 Hilbert 符号 $(a,b)_p$ 完全显式化. 回顾: $p$ 在 $K$ 中全分歧, 且 $(p)=(\pi)^{p-1}$, 其中 $\pi=1-\zeta$. 设 $K_\pi$ 表示 $K$ 在 $(\pi)$ 处的完备化, 并设 $U_i$ 表示 $K_\pi$ 中同余于 $1 \bmod \pi^{i}$ 的单位群. 我们有滤链
\[
\mathcal{O}_{K_{\pi}}^{\times}\supset U_{1}\supset U_{2}\supset\cdots\supset U_{p+1}\supset\cdots.
\]
若 $u\in U_{p+1}$, 则 $u$ 是 $K_\pi$ 中的 $p$ 次幂 (见 (VII, 注 \ref{rem:vii.A.6}(a))). 由此可以推出: $K_\pi^{\times}/K_\pi^{\times p}$ 由元素
\[
\pi,\ \zeta,\ 1-\pi^{2},\ \ldots,\ 1-\pi^{p}
\]
自由生成 (作为 $\F_p$-向量空间).

设 $\eta_i=1-\pi^{i}$, $i\ge 1$ (例如 $\eta_1=\zeta$).

\begin{proposition}\label{prop:viii.5.12}
Hilbert 配对
\[
a,b\mapsto(a,b)_{\pi}:K_{\pi}^{\times}\times K_{\pi}^{\times}\to\mu_{p}
\]
是满足下列条件的唯一的斜对称配对:

(a) $(\eta_{i},\eta_{j})_{\pi}=(\eta_{i},\eta_{i+j})_{\pi}\,(\eta_{i+j},\eta_{j})_{\pi}\,(\eta_{i+j},\pi)_{\pi}^{-j}$ 对所有 $i,j\ge 1$;

(b) $(\eta_{i},\pi)_{\pi}=
\begin{cases}
1 & \text{若 } 1\le i\le p-1,\\
\zeta & \text{若 } i=p;
\end{cases}$

(c) $(\cdot,\cdot)_{\pi}=1$ 在 $U_{i}\times U_{j}$ 上, 若 $i+j\ge p+1$.
\end{proposition}

提示见 \cite{CasselsFrohlich1967}, p.~354.

\begin{example}\label{ex:viii.5.13}
(三次互反律; Eisenstein). 设 $p=3$, 此时 $K=\Q[\zeta]$, $\zeta=\frac{-1+\sqrt{-3}}{2}$, $\pi=-\zeta\sqrt{-3}$. 则 $\mathcal{O}_K=\Z[\zeta]$, 且 $\mathcal{O}_K$ 的每个非零元素都可写成 $\zeta^{i}\pi^{j}a$ 的形式, 其中 $a\equiv\pm1 \bmod 3\mathcal{O}_K$. 此时互反律成为:
\[
\left(\frac{a}{b}\right)=\left(\frac{b}{a}\right)
\]
若 $a$ 与 $b$ 互素且同余于 $\pm 1 \bmod 3\mathcal{O}_K$; 而
\[
\begin{cases}
\left(\dfrac{\zeta}{a}\right)=\zeta^{-m-n},\\[2ex]
\left(\dfrac{\pi}{a}\right)=\zeta^{m},
\end{cases}
\]
若 $a=\pm\big(1+3(m+n\zeta)\big)$.

注意, 若 $a\in\Z$, 则 $a\equiv\pm1 \bmod 3\mathcal{O}_K$ 自动成立.
\end{example}

\subsection*{应用}

固定奇素数 $p$ 与本原 $p$ 次单位根 $\zeta$. 若整数 $x,y,z$ 满足 $x^{p}+y^{p}=z^{p}$, 则
\[
\prod_{i=0}^{p-1}(x+\zeta^{i}y)=z^{p}.
\]
不妨设 $x,y,z$ 无公因子. 若 $p\nmid xyz$, 则 $\Z[\zeta]$ 中的元素 $x+\zeta^{i}y$ 两两互素 (ANT, 6.9). 因此, 每一个生成的理想都是 $p$ 次幂, 对
\[
\alpha=\frac{x+\zeta y}{x+y}=1-\frac{y\pi}{x+y},\qquad \pi=1-\zeta
\]
亦然.

于是\footnote{此脚注修正了上一版本中一处令人困惑的排印错误; 见 \sx{667386}.}
$\left(\frac{\beta}{\alpha}\right)=1$ 对所有与 $\alpha$ 互素的 $\beta\in\Z[\zeta]$ 成立.

\begin{theorem}\label{thm:viii.5.14}
设 $x,y,z$ 是两两互素的正整数, 满足 $p\nmid xyz$ 且 $x^{p}+y^{p}=z^{p}$. 则对每个整除 $xyz$ 的素数 $q$, 有 $q^{p-1}\equiv 1 \bmod p^{2}$.
\end{theorem}

\begin{proof}
此时幂互反律成为
\[
\left(\frac{\beta}{\alpha}\right)\left(\frac{\alpha}{\beta}\right)^{-1}=\zeta^{\mathrm{Tr}_{\Q[\zeta]/\Q}(\eta)},
\]
其中 $\eta=\dfrac{\beta-1}{p}\cdot\dfrac{\alpha-1}{\pi}$. 我们把此式应用于 $\beta=q^{p-1}$. 不失一般性, 设 $q|y$, 于是 $\alpha\equiv 1 \bmod q$ 且 $\left(\frac{\alpha}{q}\right)=1$. 再者,
\[
\mathrm{Tr}(\eta)=\frac{q^{p-1}-1}{p}\,\mathrm{Tr}\Big(\frac{\alpha-1}{\pi}\Big),
\]
但
\[
\mathrm{Tr}\frac{\alpha-1}{\pi}=\mathrm{Tr}\Big(-\frac{y}{x+y}\Big)=-\frac{y}{x+y}(p-1)
\]
不被 $p$ 整除. 故 $\frac{q^{p-1}-1}{p}$ 被 $p$ 整除.
\end{proof}

\begin{corollary}[WIEFERICH 条件]\label{cor:viii.5.15}
若 $X^{p}+Y^{p}=Z^{p}$ 有解 $x,y,z$, 其中 $x,y,z$ 都是不被 $p$ 整除的正整数, 则 $2^{p-1}\equiv 1 \bmod p^{2}$.
\end{corollary}

\begin{proof}
若 $x^{p}+y^{p}=z^{p}$, 则 $x,y,z$ 中至少有一个是偶数.
\end{proof}

类似的论证 (用不同的 $\beta$) 证明了 Mirimanoff 条件: $3^{p-1}\not\equiv 1 \bmod p^{2}$.

满足 Wieferich 条件的 $<3\times 10^{9}$ 的素数只有 $1093$ 与 $3511$, 而它们都不满足 Mirimanoff 条件. 因此这证明了 Fermat 大定理第一情形对 $p<3\times 10^{9}$ 成立.

\begin{exercise}\label{exc:viii.5.16}
设 $p\in\Z$ 是同余于 $1$ 模 $3$ 的素数 (于是 $\F_p$ 包含 $1$ 的三次根). 证明: $2$ 模 $p$ 是立方当且仅当 $p$ 具有形式 $x^{2}+27y^{2}$, $x,y\in\Z$.
\end{exercise}

\begin{NOTES}
定理\ref{thm:viii.5.14} 由 Furtw\"angler 证明 (见 Hasse 1970, II, 22); 另见 \cite{Koch1992}, II, 6.3 与 Herbrand 1936, p.~47. 类域论还可以简化 Kummer 第二判别法在 Fermat 定理第二情形成立时的证明 (J.~Herbrand, Sur les classes des corps circulaires, \emph{J.~Math.~Pures Appl.}, IX.~S\'er.~11, 417--441).
\end{NOTES}
\section{数域上二次型的分类}\label{sec:VIII.6}

前面我们证明了: 数域上的非退化二次型在域中表示零当且仅当它在域的每个完备化中表示零. 本节对数域上的二次型给出完全的分类. 具体地, 我们将:

(a) 证明数域 $K$ 上的两个二次型等价当且仅当它们在 $K$ 的每个完备化上都等价;

(b) 给出局部域上二次型的完备不变量列表;

(c) 确定哪些局部不变量族来自一个整体二次型.

\subsection*{二次型的一般知识}

本小节中, $k$ 是特征 $\ne 2$ 的任意域. 设 $(V,Q)$ 是 $k$ 上带伴随双线性型 $B$ 的二次空间,
\[
B(v,w)=\frac{1}{2}\big(Q(v+w)-Q(v)-Q(w)\big),
\]
并设 $U_1,U_2$ 是 $V$ 的子空间. 若 $V$ 的每个元素都可以唯一地写成 $v=u_1+u_2$ 的形式, 其中 $u_1\in U_1$, $u_2\in U_2$, 则记 $V=U_1\oplus U_2$. 若再有 $B(u_1,u_2)=0$ 对所有 $u_1\in U_1$, $u_2\in U_2$ 成立, 则记 $V=U_1\perp U_2$. 对 $V$ 的每个子空间 $U$,
\[
U^{\perp}=\{v\in V\mid B(u,v)=0 \text{ 对所有 } u\in U\}.
\]
若 $Q|U$ 非退化, 则 $V=U\perp U^{\perp}$.

设 $(V,Q)$ 与 $(V',Q')$ 是 $k$ 上的二次空间. 态射 $s\colon(V,Q)\to(V',Q')$ 是 $k$-向量空间的线性映射 $s\colon V\to V'$, 使得 $Q'(s(v))=Q(v)$ 对所有 $v\in V$ 成立. 态射称为\textbf{同构} (isomorphism), 若它有同为态射的逆. 同构 $(V,Q)\to(V',Q')$ 也称为\textbf{等距} (isometry).

\begin{proposition}\label{prop:viii.6.1}
设 $(V,Q)$ 是二次空间. 若 $Q$ 表示 $a\in k^{\times}$, 则存在 $e\in V$ 使 $Q(e)=a$, 且存在子空间 $U\subset V$ 使 $V=U\perp k\cdot e$.
\end{proposition}

\begin{proof}
因为 $Q$ 表示 $a$, 存在 $e\in V$ 使 $Q(e)=a$, 取 $U$ 为 $k\cdot e$ 的正交补即可.
\end{proof}

设 $(V,Q)$ 是二次空间. 对每个满足 $Q(y)\ne 0$ 的 $y\in V$, 定义关于 $y$ (或关于直线 $k\cdot y$) 的对称 (symmetry) 为映射
\[
\tau_{y}(x)=x-\frac{2B(x,y)}{Q(y)}\,y.
\]
注意 $\tau_y$ 是态射 $(V,Q)\to(V,Q)$ 且 $\tau_y\circ\tau_y=\mathrm{id}$, 故 $\tau_y$ 是等距. 它把直线 $k\cdot y$ 中的每个向量变号, 并固定超平面 $(k\cdot y)^{\perp}$ 中的每个向量. 因此它是关于超平面 $(k\cdot y)^{\perp}$ 的反射.

\begin{proposition}\label{prop:viii.6.2}
设 $U$ 与 $W$ 是二次空间 $(V,Q)$ 的等距子空间, 并设 $Q|U$ 非退化. 则 $U^{\perp}$ 与 $W^{\perp}$ 等距.
\end{proposition}

\begin{proof}
对 $\dim U$ 作归纳证明. 先设 $U$ 与 $W$ 都是直线, 即 $U=ku$, $W=kw$. 则 $Q(u)\ne 0$, $Q(w)\ne 0$, 且不妨设 $Q(u)=Q(w)$. 由
\[
Q(u+w)+Q(u-w)=2Q(u)+2Q(w)=4Q(u)
\]
可见 $Q(u+w)$ 与 $Q(u-w)$ 中至少有一个非零, 必要时把 $w$ 换成 $-w$, 可设是后者. 于是对称 $\tau_{u-w}$ 有定义:
\[
\tau_{u-w}x=x-\frac{2B(x,u-w)}{Q(u-w)}\,(u-w).
\]
则 $\tau_{u-w}(u)=w$, 因为
\[
Q(u-w)=Q(u)+Q(w)-2B(u,w)=2Q(u)-2B(u,w)=2B(u,u-w),
\]
故 $\tau_{u-w}$ 把 $U^{\perp}$ 等距地映到 $W^{\perp}$.

于是可设 $\dim U\ge 2$, 从而 $U$ 有非平凡分解 $U=U_1\perp U_2$. 由于 $W\approx U$, 存在分解 $W=W_1\perp W_2$ 使 $U_1\approx W_1$, $U_2\approx W_2$. 注意 $Q|U_1$ 非退化, 且 $U_1^{\perp}=U_2\perp U^{\perp}$. 归纳假设蕴含 $U_2\perp U^{\perp}$ 等距于 $W_2\perp W^{\perp}$, 而等距的选取定义了分解 $U_2\perp U^{\perp}=X\perp Y$, 其中 $X\approx W_2$, $Y\approx W^{\perp}$. 但 $X\approx U_2$, 再由归纳假设得 $Y\approx U^{\perp}$. 故 $W^{\perp}\approx U^{\perp}$.
\end{proof}

对二次空间 $(V,Q)$ 选一组基 $e_i$, 便得到二次型
\[
q(X_{1},\ldots,X_{n})=\sum a_{ij}X_{i}X_{j},\qquad a_{ij}=B(e_{i},e_{j}).
\]
反之, 二次型 $q$ 定义二次空间 $(k^{n},q)$.

两个二次型 $q$ 与 $q'$ 称为\textbf{等价的} (equivalent), 记作 $q\sim q'$, 若它们定义同构的二次空间, 即其中一个可以由另一个经可逆变量代换得到. 若 $q$ 与 $q'$ 是两组不同变量的二次型, 则记 $q+q'$ 为 $q\perp q'$; 此时 $(k^{m+n},q\perp q')=(k^{m},q)\perp(k^{n},q')$.

由命题\ref{prop:viii.6.2} 得到:

设 $q=r\perp s$ 与 $q'=r'\perp s'$ 是两个二次型, 并设 $r$ 非退化. 若 $q\sim q'$ 且 $r\sim r'$, 则 $s\sim s'$.

由命题\ref{prop:viii.6.1} 得到:

$n$ 元非退化二次型 $q$ 表示 $a$ 当且仅当 $q\sim r\perp aZ^{2}$, 其中 $r$ 是 $n-1$ 元二次型.

二次空间 $(V,Q)$ 的\textbf{秩} (rank) 定义为矩阵 $(B(e_i,e_j))$ 对某组基 $e_i$ 的秩. 二次型 $q$ 的秩是相应二次空间的秩. 当 $q$ 写成对角形 $q=a_1X_1^{2}+\cdots+a_rX_r^{2}$ 时, $q$ 的秩是非零系数 $a_i$ 的个数, 即 $q$ 中实际出现的变量个数.

\subsection*{局部--整体原则}

\begin{theorem}[HASSE--MINKOWSKI]\label{thm:viii.6.3}
设 $q$ 与 $q'$ 是数域 $K$ 上的二次型. 若 $q$ 与 $q'$ 在所有素 $v$ 的 $K_v$ 上都等价, 则 $q$ 与 $q'$ 在 $K$ 上等价.
\end{theorem}

\begin{proof}
不妨设 $q$ 与 $q'$ 非退化. 对 $q$ 与 $q'$ 的公共秩 $n$ 作归纳. 若 $n=0$, 两个型都是零型, 无需证明. 否则, 存在 $q$ 表示的 $a\in K^{\times}$. 则 $q(X_1,\ldots,X_n)-aZ^{2}$ 在 $K$ 中表示 $0$, 从而对所有 $v$ 在 $K_v$ 中表示 $0$. 把定理\ref{thm:viii.3.5} 应用于 $q'-aZ^{2}$, 得知 $q'$ 在 $K$ 中表示 $a$. 因此 $q\sim q_1\perp aZ^{2}$, $q'\sim q_1'\perp aZ^{2}$, 其中 $q_1,q_1'$ 是 $n-1$ 秩二次型. 现在 (\ref{prop:viii.6.2}) 表明 $q_1$ 与 $q_1'$ 对所有 $v$ 在 $K_v$ 上等价, 故 (由归纳) 它们在 $K$ 上等价. 这蕴含 $q$ 与 $q'$ 在 $K$ 上等价.
\end{proof}

\begin{remark}\label{rem:viii.6.4}
设 $(V,Q)$ 是域 $k$ 上的二次空间, $O$ 是其等距群. 定理\ref{thm:viii.6.3} 说的是
\[
H^{1}(K,O)\to\prod_{v}H^{1}(K_{v},O)
\]
是单射.
\end{remark}

\subsection*{局部域上二次型的分类}

阿基米德情形. $\CC$ 上的任何二次型 (对任何特征 $\ne 2$ 的代数闭域亦然) 等价于唯一的二次型
\[
X_{1}^{2}+\cdots+X_{n}^{2}.
\]
因此 $\CC$ 上两个二次型等价当且仅当它们有相同的秩 $n$.

根据 Sylvester 定理, $\R$ 上的二次型 $q$ 等价于唯一的二次型
\[
X_{1}^{2}+\cdots+X_{r}^{2}-X_{r+1}^{2}-\cdots-X_{r+t}^{2}.
\]
$-1$ 的个数 $t$ 称为\textbf{负惯性指数} (index of negativity). 因此 $\R$ 上两个二次型等价当且仅当它们有相同的秩 $n$ 与相同的负惯性指数 $t$.

非阿基米德情形. 设 $K$ 是局部域. 回顾: Hilbert 符号 $(\cdot,\cdot)$ 对 $a,b\in K^{\times}$ 可定义为
\[
(a,b)=
\begin{cases}
1 & \Longleftrightarrow X^{2}-aY^{2}-bZ^{2} \text{ 表示 } 0\\
-1 & \text{否则,}
\end{cases}
\Longleftrightarrow aY^{2}+bZ^{2} \text{ 表示 } 1.
\]

\begin{lemma}\label{lem:viii.6.5}
Hilbert 符号具有如下性质:

(a) 它是双乘性的, 且 $(ac^{2},bd^{2})=(a,b)$ 对所有 $a,b,c,d\in K^{\times}$;

(b) 对任何非平方 $a\in K^{\times}$, 存在 $b\in K^{\times}$ 使 $(a,b)=-1$;

(c) $(b,a)=(a,b)^{-1}=(a,b)$;

(d) $(a,-a)=(1,a)=1$.
\end{lemma}

\begin{proof}
显然, 把 $a$ 或 $b$ 乘以平方不改变 $(a,b)$; (c) 也显然.

注意 $(a,b)=1$ 当且仅当 $b$ 是来自 $K[\sqrt{a}]$ 的范数. 由局部类域论, 若 $a$ 不是 $K$ 中的平方, 则 $\Nm(K[\sqrt{a}]^{\times})$ 是 $K^{\times}$ 中指数 $2$ 的子群, 故 $b\mapsto(a,b)$ 是同构 $K^{\times}/\Nm(K[\sqrt{a}]^{\times})\to\{\pm 1\}$. 这完成了 (a) 与 (b) 的证明. 最后, $aX^{2}-aY^{2}=a(X^{2}-Y^{2})$, 而 $X^{2}-Y^{2}$ 表示 $a^{-1}$, 因为它表示 $0$.
\end{proof}

若 $q\sim a_{1}X_{1}^{2}+\cdots+a_{n}X_{n}^{2}$, 其中 $a_{1},\ldots,a_{n}\in K^{\times}$, 则令
\[
\begin{aligned}
n(q)&=n,\\
d(q)&=a_{1}\cdots a_{n}\qquad(\text{属于 } K^{\times}/K^{\times 2}),\\
S(q)&=\prod_{1\le i\le j\le n}(a_{i},a_{j})=\prod_{1\le i\le n}(a_{i},d_{i})\qquad(\text{属于 } \{\pm 1\}),
\end{aligned}
\]
其中 $d_i=a_1\ldots a_i$. 于是 $n(q)$ 是 $q$ 的秩, $d(q)$ 是 $q$ 的判别式, 两者都只依赖于 $q$ 的等价类. 我们将证明 $S(q)$ 亦然. 它称为 $q$ 的 \textbf{Hasse 不变量}\index{hassebubianliang@Hasse 不变量(Hasse invariant)}.

\begin{remark}\label{rem:viii.6.6}
\cite{Serre1970} 定义
\[
\varepsilon(q)=\prod_{1\le i<j\le n}(a_{i},a_{j})=\prod(a_{i},d_{i-1}).
\]
注意
\[
S(q)=\varepsilon(q)\prod_{i=1}^{n}(a_{i},a_{i})=\varepsilon(q)\prod_{i=1}^{n}(-1,a_{i})(-a_{i},a_{i})=\varepsilon(q)(-1,d(q)).
\]
因此知道 $(d(q),S(q))$ 等价于知道 $(d(q),\varepsilon(q))$.
\end{remark}

\begin{proposition}\label{prop:viii.6.7}
元素 $S(q)$ 只依赖于 $q$ 的等价类.
\end{proposition}

\begin{proof}
只需证明 $\varepsilon(q)$ 只依赖于 $q$ 的等价类. 当 $q$ 的秩为 $1$ 时无需证明: 对所有 $q$, $\varepsilon(q)=1$ (空积).

其次设 $q\sim aX^{2}+bY^{2}\sim a'X^{2}+b'Y^{2}$. 由于二者等价, $aX^{2}+bY^{2}$ 与 $a'X^{2}+b'Y^{2}$ 要么都表示 $1$, 要么都不表示 $1$, 故 $(a,b)=(a',b')$.

再设 $n>2$ 且
\[
q\sim a_{1}X_{1}^{2}+\cdots+a_{i}X_{i}^{2}+a_{i+1}X_{i+1}^{2}+\cdots\sim a_{1}^{\prime}X_{1}^{2}+\cdots+a_{i}^{\prime}X_{i}^{2}+a_{i+1}^{\prime}X_{i+1}^{2}+\cdots
\]
其中 $a_{j}=a_{j}^{\prime}$ 对 $j\ne i,i+1$ 成立. 我们要证明
\[
(a_{i},d_{i-1})(a_{i+1},d_{i})=(a_{i}^{\prime},d_{i-1}^{\prime})(a_{i+1}^{\prime},d_{i}^{\prime}).
\]
但
\[
(a_{i},d_{i-1})(a_{i+1},d_{i})=(a_{i},d_{i-1})(a_{i+1},d_{i-1})(a_{i+1},a_{i})=(a_{i}a_{i+1},d_{i-1})(a_{i},a_{i+1}),
\]
而 $a_{i}a_{i+1}$ 与 $a_{i}^{\prime}a_{i+1}^{\prime}$ 相差一个平方, 故剩下只需证 $(a_{i},a_{i+1})=(a_{i}^{\prime},a_{i+1}^{\prime})$. 根据命题\ref{prop:viii.6.2}, $a_{i}X_{i}^{2}+a_{i+1}X_{i+1}^{2}\sim a_{i}^{\prime}X_{i}^{2}+a_{i+1}^{\prime}X_{i+1}^{2}$, 而前面已证明这蕴含 $(a_{i},a_{i+1})=(a_{i}^{\prime},a_{i+1}^{\prime})$.
\end{proof}

下面的初等引理完成了证明.

\begin{lemma}\label{lem:viii.6.8}
设 $B$ 与 $B'$ 是非退化二次空间 $(V,Q)$ 的正交基. 则存在一列正交基 $B_1,B_2,\ldots,B_m$ 使 $B_1=B$, $B_m=B'$, 且每个 $B_i$ 都由 $B_{i-1}$ 改动至多两个相邻元素得到.
\end{lemma}

\begin{proof}
见 \cite{OMeara1963}, 引理 58.1.
\end{proof}

\begin{proposition}\label{prop:viii.6.9}
设 $q$ 是非阿基米德局部域 $K$ 上的 $n$ 元非退化二次型, $a\in K^{\times}$. 则 $q$ 表示 $a$ 当且仅当

(a) $n=1$ 且 $a=d(q)$ (在 $K^{\times}/K^{\times 2}$ 中);

(b) $n=2$ 且 $(a,-d)(-1,d)=S(q)$ (等价地, $(a,-d)=\varepsilon(q)$);

(c) $n=3$ 且要么 $a\ne -d(q)$ (模平方), 要么 $a=-d(q)$ (模平方) 且 $(-1,-1)=S(q)$;

(d) $n\ge 4$.
\end{proposition}

\begin{proof}
(a) 显然 $dX^{2}$ 表示 $a$ 当且仅当 $a=d$ (在 $K^{\times}/K^{\times 2}$ 中).

(b) 设 $q=bX^{2}+cY^{2}$. 显然 $bX^{2}+cY^{2}$ 表示 $a$ 当且仅当 $abX^{2}+acY^{2}$ 表示 $1$, 即 $(ab,ac)=1$. 但
\[
(ab,ac)=(a,a)(a,b)(a,c)(b,c)=(a,-1)(a,d(q))(b,c)=(a,-d(q))\cdot\varepsilon(q),
\]
故条件为
\[
\varepsilon(q)=(a,-d(q)).
\]

(c) 设 $q=a_1X_1^{2}+a_2X_2^{2}+a_3X_3^{2}$. 则 $q$ 表示 $a$ 当且仅当存在 $e\in K^{\times}$ 使方程
\[
a_{1}X_{1}^{2}+a_{2}X_{2}^{2}=e=a_{3}X_{3}^{2}-aX_{4}^{2}
\]
有解. 根据 (b), 这当且仅当
\begin{equation*}
(e,-a_{1}a_{2})=(a_{1},a_{2}),\qquad (e,a_{3}a)=(a_{3},-a)
\tag{$*$}
\end{equation*}
时成立. 考虑 $\F_2$-向量空间 $V$ ($\dim V\ge 2$) 上的两个线性型 $f,g\colon V\to\F_2$. 联立线性方程 $f(x)=\varepsilon_1$, $g(x)=\varepsilon_2$ 有解, 除非它们不相容, 即: $f=0$ 且 $\varepsilon_1=-1$; 或 $g=0$ 且 $\varepsilon_2=-1$; 或 $f=g$ 且 $\varepsilon_1=-\varepsilon_2$.

把这个观察应用于线性型 $(\cdot,-a_1a_2)$, $(\cdot,a_3a)\colon K^{\times}/K^{\times 2}\to\{\pm 1\}$, 我们发现: 存在满足 (*) 的 $e$, 除非 $-a_1a_2=a_3a$ (在 $K^{\times}/K^{\times 2}$ 中) 且 $(a_1,a_2)=-(a_3,a)$. 第一个等式说 $a=-d(q)$ (模平方), 而 (当 $a=-d(q)$ 时) 第二个说 $(-1,-1)=S(q)$.

(d) 此时 $q(X_1,\ldots,X_n)-aZ^{2}$ 的秩 $\ge 5$, 故它表示 $0$ (见 \ref{cor:viii.3.10}).
\end{proof}

\begin{theorem}\label{thm:viii.6.10}
非阿基米德局部域上的两个二次型等价当且仅当它们有相同的秩、相同的判别式与相同的 Hasse 不变量.
\end{theorem}

\begin{proof}
我们在命题\ref{prop:viii.6.7} 中证明了等价的型有相同的不变量. 反之, 对两个型 $q$ 与 $q'$ 的公共秩 $n$ 作归纳. 秩为 $1$ 的两个二次型显然在判别式相同时等价, 故可设 $n>1$. 由命题\ref{prop:viii.6.9}, $q$ 与 $q'$ 表示 $K^{\times}$ 中相同的元素. 特别地, 存在同时被 $q$ 与 $q'$ 表示的 $a\in K^{\times}$. 于是
\[
q\sim q_{1}\perp aZ^{2},\qquad q'\sim q_{1}^{\prime}\perp aZ^{2},
\]
其中 $q_1,q_1'$ 是 $n-1$ 秩二次型. 现在
\[
d(q)=a\cdot d(q_{1}),\qquad S(q)=(a,d(q_{1}))\cdot S(q_{1}),
\]
对 $q'$ 与 $q_1'$ 亦然. 因此 $q_1$ 与 $q_1'$ 有相同的不变量, 归纳假设表明 $q_1\sim q_1'$.
\end{proof}

\begin{proposition}\label{prop:viii.6.11}
设 $q$ 是非阿基米德局部域 $K$ 上秩为 $n$ 的二次型.

(a) 若 $n=1$, 则 $S(q)=(-1,d)$.

(b) 若 $n=2$, 则 $d(q)=-1$ (模平方) 蕴含 $S(q)=(-1,-1)$.

除这些限制之外, 每个三元组 $n\ge 1$, $d\in K^{\times}/K^{\times 2}$, $s=\pm 1$ 都可以作为 $K$ 上某个二次型的不变量组出现.
\end{proposition}

\begin{proof}
情形 $n=1$. 此时 $q=dX^{2}$, 而 $S(q)=(d,d)=(-1,d)$.

情形 $n=2$. 对 $q=aX^{2}+bY^{2}$, $S(q)=(a,a)(b,d)$, 故
\[
d=-1\Rightarrow S(q)=(-1,a)(-1,b)=(-1,d)=(-1,-1).
\]
反之, 型 $X^{2}-Y^{2}$ 满足 $d=-1$, $S=(-1,-1)$.

现在设 $d\ne -1$ 与 $s$ 已给定. 我们找 $a\in K^{\times}$ 使 $q\overset{\text{def}}{=}aX^{2}+adY^{2}$ 满足 $S(q)=s$. 但
\[
S(q)=(a,a)(ad,d)=(a,-1)(a,d)(d,d)=(a,-d)(d,d).
\]
由于 $-d\ne 1$ (在 $K^{\times}/K^{\times 2}$ 中), 可选 $a$ 使 $(a,-d)=s\cdot(d,d)$.

情形 $n=3$. 取 $a\in K^{\times}$ 使 $a\ne -d$ (在 $K^{\times}/K^{\times 2}$ 中). 由于对 $a$ 的这个条件, 存在秩为 $2$ 的二次型 $q_1$ 使
\[
d=d(q_{1})a,\qquad s=S(q_{1})(a,d).
\]
取 $q=q_1\perp aZ^{2}$.

情形 $n\ge 4$. 存在形状为
\[
q_{1}(X_{1},X_{2},X_{3})+X_{4}^{2}+\cdots+X_{n}^{2}
\]
且具有所需不变量的二次型.
\end{proof}

\paragraph{一般说明}
对 $\R$ 或 $\CC$ 上的二次型 $q$, 用与非阿基米德情形相同的公式定义 Hasse 不变量. 对 $\CC$, 恒有 $S(q)=1$; 对 $\R$, $S(q)=(-1)^{t(t+1)/2}$, 其中 $t$ 是负惯性指数 (因为 $(-1,-1)=-1$). 注意在第二种情形 $d(q)=(-1)^{t}$ (在 $\R^{\times}/\R^{\times 2}$ 中), 且当 $r\le 3$ 时 $d(q)$ 与 $S(q)$ 确定 $t$, 而当 $r>3$ 时不能.

我们说系统 $(n,d,s,\ldots)$ ($n\in\mathbb{N}$, $d\in K^{\times}/K^{\times 2}$, $s\in\{\pm1\}$, \ldots) 是\textbf{可实现的} (realizable), 若存在二次型 $q$ 使 $n(q)=n$, $d(q)=d$, $S(q)=s$, \ldots

(a) 对非阿基米德局部域 $K$, $(n,d,s)$ 可实现, 只要 $n=1$ 时 $s=1$, 且 $n=2$、$d=-1$ 时 $s=(-1,-1)$.

(b) 对 $\R$, $(n,d,s,t)$ 可实现, 只要 $0\le t\le n$, $d=(-1)^{t}$, $s=(-1)^{t(t+1)/2}$.

\subsection*{整体域上二次型的分类}

\begin{theorem}\label{thm:viii.6.12}
设 $n\in\mathbb{N}$, 并设对数域 $K$ 的每个素 $v$ 给定了 $K_v$ 上秩为 $n$ 的非退化二次型 $q(v)$. 则存在 $K$ 上的二次型 $q_0$ 使 $(q_0)_v\sim q_v$ 对每个 $v$ 成立, 当且仅当

(a) 存在 $d_0\in K^{\times}$ 使 $d_0\equiv d(q_v) \bmod K_v^{\times 2}$ 对所有 $v$;

(b) $S(q(v))=1$ 对几乎所有 $v$, 且 $\prod_{v}S(q(v))=1$.
\end{theorem}

这些条件显然是必要的. 鉴于命题\ref{prop:viii.6.11} 及其后的说明, 可以把该定理重述如下. 设给定:

$\diamond$ 一个 $n\ge 1$ 与一个 $d_0\in K^{\times}/K^{\times 2}$;

$\diamond$ 对每个素 $v$ (有限的或实的), 一个 $s_v\in\{\pm 1\}$;

$\diamond$ 对每个实素 $v$, 一个整数 $t_v$.

则: 存在 $K$ 上秩为 $n$、判别式为 $d_0$、Hasse 不变量 $S_v(q)=s_v$ 对所有 $v$、负惯性指数 $t_v(q)=t_v$ 对所有实素 $v$ 的二次型 $q$, 当且仅当

(a) $s_v=1$ 对除有限多个之外的所有 $v$, 且 $\prod_v s_v=1$;

(b) 若 $n=1$, 则 $s_v=(-1,d)_v$; 若 $n=2$, 则要么 $d\ne -1$ (在 $K_v^{\times}/K_v^{\times 2}$ 中), 要么 $s_v=(-1,-1)_v$;

(c) 对所有实素 $v$, $0\le t_v\le n$, $d_v=(-1)^{t_v}$ (模平方), 且 $s_v=(-1)^{t_v(t_v+1)/2}$.

在 $n=1$ 的情形, $q_v=d(q_v)X^{2}$, 可取 $q_0=d_0X^{2}$, 其中 $d_0$ 是由 (a) 保证存在的 $K^{\times}$ 中的元素.

关键情形是 $n=2$, 为此需要下面的引理, 其证明需要类域论.

\begin{lemma}\label{lem:viii.6.13}
设 $T$ 是 $K$ 的由实素或有限素组成的有限集, $b\in K^{\times}$. 若 $T$ 有偶数个元素, 且 $b$ 对任何 $v\in T$ 都不在 $K_v^{\times}$ 中成为平方, 则存在 $a\in K^{\times}$ 使
\[
(a,b)_{v}=
\begin{cases}
-1 & \text{对 } v\in T,\\
1 & \text{否则.}
\end{cases}
\]
\end{lemma}

\begin{proof}
(仿照 Tate, 1976, 5.2). 设 $L$ 是 $K$ 的所有指数为 $2$ 的 Abel 扩张的复合, 并令 $G=\Gal(L/K)$. 由类域论,
\[
G\simeq C_{K}/2C_{K}\simeq\mathbb{I}/K^{\times}\cdot\mathbb{I}^{2}.
\]
列
\[
0\to\mu_{2}\to L^{\times}\xrightarrow{\ x\mapsto x^{2}\ }L^{\times 2}\to0
\]
的上同调列是正合列
\[
K^{\times}\xrightarrow{\ x\mapsto x^{2}\ }K^{\times}\cap L^{\times 2}\to\Hom_{\mathrm{conts}}(G,\mu_{2})\to0.
\]
$K^{\times}$ 的每个元素在 $L$ 中都成为平方, 故有同构
\[
K^{\times}/K^{\times 2}\to\Hom_{\mathrm{conts}}(\mathbb{I}/K^{\times}\cdot\mathbb{I}^{2},\mu_{2}).
\]
该映射把 $a\in K^{\times}$ 送到连续同态
\[
(c_{v})\mapsto\prod(a,c_{v})_{v}
\]
(因为 Hilbert 符号与局部 Artin 映射之间的关系). 于是, 找 $a$ 等价于找同态 $f\colon\mathbb{I}/\mathbb{I}^{2}\to\mu_2$ 使得

(a) $f=1$ 在主 idèle 上;

(b) $f(1,\ldots,1,i_{v}(b),1,\ldots,1)=
\begin{cases}
-1 & \text{对 } v\in T,\\
1 & \text{否则,}
\end{cases}$

其中 $i_v$ 是包含 $K\hookrightarrow K_v$. 对每个 $v$, 设 $B_v$ 是由 $i_v(b)$ 生成的 $K_v^{\times}/K_v^{\times 2}$ 的 $\F_2$-子空间, 并设 $B=\prod B_v\subset\mathbb{I}/\mathbb{I}^{2}$. 由于 $i_v(b)$ 对 $v\in T$ 不是平方, 存在 (自动连续的) 线性型 $f_1\colon B\to\mu_2$ 满足条件 (b), 而 $f_1$ 延拓为 $\mathbb{I}/\mathbb{I}^{2}$ 上满足 (a) 的连续线性型当且仅当 $f_1$ 在 $B$ 中每个主 idèle 上取值 $1$. $f_1$ 在 $b$ 的主 idèle 上的值为 $\prod_{v\in T}(-1)$, 由 $T$ 含偶数个元素的假设它等于 $1$. 设 $c\in K^{\times}$ 使其 idèle 属于 $B$. 则对每个 $v$, $i_v(c)=1$ 或 $i_v(b)$ (在 $K_v^{\times}/K_v^{\times 2}$ 中). 因此, $i_v(c)$ 对所有 $v$ 在 $K_v[\sqrt{b}]$ 中成为平方, 而这 (由 \ref{thm:viii.1.1}) 蕴含 $c$ 是 $K[\sqrt{b}]$ 中的平方. 故在 $K^{\times}/K^{\times 2}$ 中 $c=1$ 或 $b$, 从而 $f_1$ 在其 idèle 上取值 $1$.
\end{proof}

现在证明定理的 $n=2$ 情形. 设对所有 $v$ 给定二次型 $q(v)=a(v)X^{2}+a(v)d(v)Y^{2}$, 我们寻找 $q=aX^{2}+ad_{0}Y^{2}$ 使 $q\sim q(v)$ 对所有 $v$ 在 $K_v$ 上成立. 即找 $a\in K^{\times}$ 使
\[
S_{v}(q)\overset{\text{def}}{=}(a,a)_{v}(ad_{0},d_{0})_{v}=S(q(v))
\]
对所有 $v$ 成立. 现在
\[
(a,a)_{v}(ad_{0},d_{0})_{v}=(a,-1)_{v}(a,d_{0})_{v}(-1,d_{0})_{v}=(a,-d_{0})_{v}(-1,d_{0})_{v}.
\]
把引理应用于 $T$ 为满足 $S(q(v))(-1,d_{0})_{v}=-1$ 的素之集, $b=-d_0$. 由定理的条件 (b), $T$ 是有限集, 且它有偶数个元素, 因为
\[
\prod_{v}S(q(v))(-1,d_{0})_{v}=\prod_{v}S(q(v))\cdot\prod_{v}(-1,d_{0})_{v}=1\times 1=1.
\]
再者,
\[
S(q(v))\cdot(-1,d_{0})_{v}=(a(v),-d(v))_{v},
\]
故当 $v\in T$ 时 $-i_v(d_0)=-d(v)\ne 1$ (在 $K_v^{\times}/K_v^{\times 2}$ 中). 于是引理给出了所需的元素 $a$.

接着证明 $n=3$ 情形. 对型 $q=q_1+aZ^{2}$ ($a\in K^{\times}$),
\[
n(q_{1})=n(q)-1,\qquad d(q_{1})=a\cdot d(q),\qquad S_{v}(q_{1})=(a,d(q))_{v}\cdot S_{v}(q).
\]
我们找 $a$ 使不变量 $(2,a\cdot d_v,(a,d_v)_v\cdot s_v)$ 对所有 $v$ 可实现, 即 $i_v(a)d_v=-1\Rightarrow (a,d_v)_v\cdot s_v=1$. 令 $T=\{v\mid s_v\ne 1\}$——由假设它是有限集. 由弱逼近定理 (ANT, 7.20), 存在 $a\in K^{\times}$ 使 $i_v(a)d_v\ne -1$ 对所有 $v\in T$. 此时对 $v\in T$, $d(q_1)\ne -1$, 故 $(2,d_v,s_v)$ 可实现. 对 $v\notin T$, $(a,d_v)\cdot s_v=(a,d_v)$, 而 $i_v(a)d_v=-1$ 蕴含 $(a,d_v)_v=(-d_v,d_v)_v=1$. 因此, 对这样的 $a$, 存在秩为 $2$ 的二次型 $q_1$ 使 $q_1\perp aZ^{2}$ 具有所需的不变量.

$n\ge 4$ 情形用归纳法证明. 若对所有实素 $t_v<n$, 可找到形状为 $q_1(X)\perp Z^{2}$ 且具有正确局部不变量的二次型. 若没有 $t_v=0$, 可找到形状为 $q_1(X)\perp(-Z^{2})$ 且具有正确不变量的二次型. 一般情形下, 用弱逼近定理 (ANT, 7.20) 找一个在 $t_v<n$ 的实素处为正、在 $t_v=0$ 的实素处为负的元素 $a$. 于是归纳假设使我们能找到 $q_1(X)$ 使 $q_1(X)\perp aZ^{2}$ 具有正确的不变量.

\subsection*{应用}

\begin{proposition}[GAUSS]\label{prop:viii.6.14}
正整数 $n$ 是三个平方数之和当且仅当它不具有形式 $4^{a}(8b-1)$, 其中 $a,b\in\Z$.
\end{proposition}

\begin{proof}
把上述理论应用于二次型 $X_1^{2}+X_2^{2}+X_3^{2}-nZ^{2}$——见 \cite{Serre1970}, 第 IV 章.
\end{proof}

\section{密度定理}\label{sec:VIII.7}

本节始终设 $K$ 是数域.

\begin{theorem}\label{thm:viii.7.1}
对 $K$ 的任意模 $\mathfrak{m}$ 与任意非平凡 Dirichlet 特征 $\chi\colon C_m\to\CC^{\times}$, 有 $L(1,\chi)\ne 0$.
\end{theorem}

\begin{proof}
正如第 \ref{ch:VI} 章末尾所指出的, 一旦有了互反律, 这就由定理 (VI, \ref{thm:vi.4.9}) 的证明推出.
\end{proof}

\begin{theorem}\label{thm:viii.7.2}
设 $\mathfrak{m}$ 是 $K$ 的模, $H$ 是 $\mathfrak{m}$ 的同余子群: $I^{m}\supset H\supset i(K_{m,1})$. 对任意类 $\mathfrak{k}\in I^{m}/H$, $\mathfrak{k}$ 中的素理想集具有 Dirichlet 密度 $1/(I^{m}:H)$.
\end{theorem}

\begin{proof}
把定理\ref{thm:viii.7.1} 与定理 (VI, \ref{thm:vi.4.8}) 结合.
\end{proof}

\begin{corollary}\label{cor:viii.7.3}
设 $L/K$ 是以 $G$ 为 Galois 群的有限 Abel 扩张, $\sigma\in G$. 则在 $L$ 中不分歧且满足 $(\mathfrak{p},L/K)=\sigma$ 的 $K$ 的素理想 $\mathfrak{p}$ 之集具有 Dirichlet 密度 $\frac{1}{[L:K]}$.
\end{corollary}

\begin{proof}
互反律 (V, \ref{thm:v.3.5}) 说: 对某个模 $\mathfrak{m}$ 与某个 $H\supset i(K_{m,1})$, Artin 映射定义同构 $I^{m}/H\to\Gal(L/K)$. 把定理应用于 $\sigma$ 在 $I^{m}/H$ 中的原像即可.
\end{proof}

\begin{theorem}[CHEBOTAREV]\label{thm:viii.7.4}
设 $L$ 是数域 $K$ 的有限 Galois 扩张, Galois 群为 $G$, 并设 $C\subset G$ 在共轭下稳定, 即
\[
x\in C,\ \tau\in G\Rightarrow\tau x\tau^{-1}\in C.
\]
令
\[
T=\{\mathfrak{p}\mid \mathfrak{p} \text{ 在 } L \text{ 中不分歧},\ (\mathfrak{p},L/K)\subset C\}.
\]
则 $T$ 具有 Dirichlet 密度
\[
\delta(T)=\frac{C \text{ 中元素的个数}}{G \text{ 中元素的个数}}.
\]
\end{theorem}

\begin{proof}
只需对 $C$ 是单个元素 $\sigma$ 的共轭类
\[
C=\{\tau\sigma\tau^{-1}\mid\tau\in G\}
\]
的情形证明. 设 $\sigma$ 的阶为 $f$, 并令 $M=L^{\langle\sigma\rangle}$. 则 $L$ 是 $M$ 的 $f$ 次循环扩张, 故 Artin 映射给出同构
\[
C_{\mathfrak{m}}/H\to\langle\sigma\rangle
\]
对 $M$ 的某个模 $\mathfrak{m}$ 与 $H=M^{\times}\cdot\Nm_{L/M}C_{L,\mathfrak{m}}$ 成立. 我们用记号
\[
\mathfrak{P}\,|\,\mathfrak{q}\,|\,\mathfrak{p}
\]
表示 $\mathcal{O}_L$ 的素 $\mathfrak{P}$、$\mathcal{O}_M$ 的素 $\mathfrak{q}$ 与 $\mathcal{O}_K$ 的素 $\mathfrak{p}$. 令 $d=[L:K]=(G:1)$, 并令 $c$ 为 $C$ 的阶. 我们要证
\[
\delta(T)=\frac{c}{d}.
\]
证明中我们忽略 (有限多个) 不与 $\mathfrak{m}$ 互素的素理想.

令
\[
T_{M,\sigma}=\{\mathfrak{q}\subset\mathcal{O}_{M}\mid(\mathfrak{q},L/M)=\sigma,\ f(\mathfrak{q}/\mathfrak{p})=1\}.
\]
Abel 扩张的 Chebotarev 密度定理 (\ref{cor:viii.7.3}) 表明: 满足 $T_{M,\sigma}$ 定义中第一个条件的素集的密度为 $\frac{1}{f}$, 由此 (见 (VI, \ref{prop:vi.4.5})) $T_{M,\sigma}$ 本身的密度为 $\frac{1}{f}$.

令
\[
T_{L,\sigma}=\{\mathfrak{P}\subset\mathcal{O}_{L}\mid(\mathfrak{P},L/K)=\sigma\}.
\]
我们将证明:

(a) 映射 $\mathfrak{P}\mapsto\mathfrak{q}=\mathfrak{P}\cap\mathcal{O}_M$ 定义双射 $T_{L,\sigma}\to T_{M,\sigma}$;

(b) 映射 $\mathfrak{P}\mapsto\mathfrak{p}=\mathfrak{P}\cap\mathcal{O}_K\colon T_{L,\sigma}\to T$ 恰好把 $\frac{d}{cf}$ 个 $T_{L,\sigma}$ 中的素送到 $T$ 的每个素.

把这两个断言结合起来: 映射 $\mathfrak{q}\mapsto\mathfrak{p}=\mathfrak{q}\cap\mathcal{O}_K$ 定义 $\frac{d}{cf}:1$ 的映射 $T_{M,\sigma}\to T$. 对这样的 $\mathfrak{q}$, $\Nm_{M/K}\mathfrak{q}=\mathfrak{p}$, 故 $\mathbb{N}\mathfrak{q}=\mathbb{N}\mathfrak{p}$; 于是
\[
\sum_{\mathfrak{p}\in T}\frac{1}{\mathbb{N}\mathfrak{p}^{s}}=\frac{cf}{d}\sum_{\mathfrak{q}\in T_{M,\sigma}}\frac{1}{\mathbb{N}\mathfrak{q}^{s}}\sim\frac{cf}{d}\cdot\frac{1}{f}\log\frac{1}{s-1}=\frac{c}{d}\log\frac{1}{s-1},
\]
如所欲证. 剩下要证 (a) 与 (b).

设 $\mathfrak{P}\in T_{L,\sigma}$, 并令 $\mathfrak{q}=\mathfrak{P}\cap\mathcal{O}_M$, $\mathfrak{p}=\mathfrak{P}\cap\mathcal{O}_K$. 则 $L_{\mathfrak{P}}/K_{\mathfrak{p}}$ 的 Galois 群由 $\sigma$ 生成, 但 $\sigma$ 固定 $M_{\mathfrak{q}}$, 故 $M_{\mathfrak{q}}=K_{\mathfrak{p}}$. 因此 $f(\mathfrak{q}/\mathfrak{p})=1$, 这表明 $\mathfrak{q}\in T_{M,\sigma}$, 于是有映射
\[
\mathfrak{P}\mapsto\mathfrak{q}\overset{\text{def}}{=}\mathfrak{P}\cap\mathcal{O}_{M}:T_{L,\sigma}\to T_{M,\sigma}.
\]
它是单射, 因为 $f(\mathfrak{P}/\mathfrak{q})=f(\mathfrak{q}/\mathfrak{p})^{-1}f(\mathfrak{P}/\mathfrak{p})=1\times f$, 故 $\mathfrak{P}$ 是 $L$ 中位于 $\mathfrak{q}$ 之上的唯一素. 它是满射, 因为对位于 $\mathfrak{q}\in T_{M,\sigma}$ 之上的任何素 $\mathfrak{P}$,
\[
(\mathfrak{P},L/K)=(\mathfrak{P},L/K)^{f(\mathfrak{q}/\mathfrak{p})}=(\mathfrak{P},L/M)=\sigma
\]
($\mathfrak{q}$ 属于 $T_{M,\sigma}$ 的第一个条件), 故 $\mathfrak{P}\in T_{L,\sigma}$. 这证明了 (a).

固定 $\mathfrak{p}_0\in T$, 并设 $\mathfrak{P}_0\in T_{L,\sigma}$ 位于 $\mathfrak{p}_0$ 之上. 则对 $\tau\in G$,
\[
(\tau\mathfrak{P}_{0},L/K)=\tau(\mathfrak{P}_{0},L/K)\tau^{-1},
\]
故
\[
\tau(\mathfrak{P}_{0},L/K)\tau^{-1}=\sigma\iff\tau\in C_{G}(\sigma)
\]
($\sigma$ 在 $G$ 中的中心化子). 因此映射 $\tau\mapsto\tau\mathfrak{P}_0$ 是双射
\[
C_{G}(\sigma)/G(\mathfrak{P}_{0})\to\{\mathfrak{P}\in T_{L,\sigma}\mid\mathfrak{P}\cap\mathcal{O}_{K}=\mathfrak{p}_{0}\}.
\]
分解群 $G(\mathfrak{P}_0)$ 等于 $\langle\sigma\rangle$, 其阶为 $f$; 而 $C_G(\sigma)$ 的阶为 $\frac{d}{c}$, 因为存在双射
\[
\tau\mapsto\tau\sigma\tau^{-1}:G/C_{G}(\sigma)\to C.
\]
因此 $(C_G(\sigma):G(\mathfrak{P}_0))=\frac{d}{cf}$, 于是我们证明了: 对每个 $\mathfrak{p}\in T$, 恰有 $\frac{d}{cf}$ 个素 $\mathfrak{P}\in T_{L,\sigma}$ 位于 $\mathfrak{p}$ 之上. 这证明了 (b), 从而完成了定理的证明.
\end{proof}

\begin{remark}\label{rem:viii.7.5}
关于 Chebotarev 密度定理的有效形式, 见 Lagarias 与 Odlyzko (Algebraic Number Fields, Ed.~Fr\"ohlich, 1977). 设 $L$ 是 $K$ 的有限 Galois 扩张, 并令
\[
\pi_{C}(x)=|\{\mathfrak{p}\mid(\mathfrak{p},L/K)\subset C,\ \mathbb{N}\mathfrak{p}\le x\}|.
\]
则
\[
\pi_{C}(x)=\frac{c}{d}\cdot\frac{x}{\log x}+\text{具体的误差项}.
\]
\end{remark}

\section{函数域; 几何类域论}\label{sec:VIII.8}

我们还应当包括函数域 (某个 $\F_p(T)$ 的有限扩张) 的类域论. 为此, 既可以模仿数域情形的证明 (见 \cite{ArtinTate1961}), 也可以 (更好) 把证明建立在 Tsen 定理之上 (见 \cite{Milne2006}, 第 I 章附录).

设 $K$ 是域 $k$ 上一个变量的函数域, 即 $k(T)$ 的有限扩张. 几何类域论描述 $K$ 的``几何''Abel 扩张, 即``不来自 $k$''的 Abel 扩张. 关于几何类域论的文献列表, 见 \mo{73054}.

\section{数域的上同调}\label{sec:VIII.9}

我们还应当包括 Nakayama 与 Tate 的定理的证明 (例如 J.~Tate, The cohomology groups of tori in finite Galois extensions of number fields, \emph{Nagoya Math.~J.}, 27, 1966, 709--719) 以及 Poitou 与 Tate 的定理的证明 (例如 \cite{Milne2006}, 第 I 章).

\section{再谈 \texorpdfstring{$L$}{L}-级数}\label{sec:VIII.10}

设 $\chi$ 是一个 Dirichlet $L$-级数对应的特征. 则存在常数 $A(\chi)\ge 0$, $a(\chi)$, $b(\chi)\in\CC$, 使得
\[
\Phi(s,\chi)\overset{\text{def}}{=}A(\chi)^{s}\,\Gamma\Big(\frac{s}{2}\Big)^{a(\chi)}\Gamma\Big(\frac{s+1}{2}\Big)^{b(\chi)}L(s,\chi)
\]
满足函数方程
\[
\Phi(s,\chi)=W(\chi)\,\Phi(1-s,\bar{\chi}),\qquad |W(\chi)|=1.
\]
见 \cite{Narkiewicz1990}.

\subsection*{Artin $L$-级数}

设 $L$ 是 $K$ 的以 $G$ 为 Galois 群的有限 Galois 扩张. 设 $V$ 是 $\CC$ 上的有限维向量空间, 并设
\[
\rho\colon G\to\mathrm{Aut}_{\CC}(V)
\]
是 $G$ 到 $V$ 的线性自同构群的同态. 我们称 $\rho$ 为 $G$ 在 $\CC$ 上的一个 (有限维) 表示. $\rho$ 的迹是映射
\[
\sigma\mapsto\chi(\sigma)=\mathrm{Tr}(\rho(\sigma)).
\]
(回顾: $m\times m$ 矩阵 $(a_{ij})$ 的迹是 $\sum a_{ii}$, 自同态的迹是其相对任一组基的矩阵的迹.) 对 $\sigma\in G$, 令
\[
P_{\sigma}(T)=\det(1-\rho(\sigma)T\mid V)=\prod_{i=1}^{\dim V}(1-a_{i}T),\qquad a_{i}\in\CC
\]
为 $\rho(\sigma)$ 的特征多项式. 注意 $P_{\sigma}(T)$ 只依赖于 $\sigma$ 的共轭类, 故对 $K$ 的每个在 $L$ 中不分歧的素 $\mathfrak{p}$, 可定义
\[
P_{\mathfrak{p}}(T)=P_{\sigma}(T),\qquad \sigma=(\mathfrak{P},L/K),\ \mathfrak{P}|\mathfrak{p} \text{ 某个}.
\]
对这样的 $\mathfrak{p}$, 令
\[
L_{\mathfrak{p}}(s,\rho)=\frac{1}{P_{\mathfrak{p}}(\mathbb{N}\mathfrak{p}^{-s})}
\]
并令
\[
L(s,\rho)=\prod L_{\mathfrak{p}}(s,\rho).
\]
例如, 若 $L/K$ 是 Abel 的, 则表示可对角化\footnote{我们的意思是: 相对于 $V$ 的某组基,
\[
\rho(g)=\begin{pmatrix}
\chi_{1}(g) & 0 & \cdots & 0\\
0 & \chi_{2}(g) & \cdots & 0\\
\vdots & \vdots & \ddots & \vdots\\
0 & 0 & \cdots & \chi_{m}(g)
\end{pmatrix},\qquad g\in G.
\]}
\[
\rho\approx\chi_{1}\oplus\cdots\oplus\chi_{m},
\]
其中每个 $\chi_i$ 都是同态 $G\to\CC^{\times}$. 与 Artin 映射
\[
C_{m}\to G
\]
复合之后, $\chi_i$ 成为 Dirichlet 特征, 于是 Artin $L$-级数等同于 Dirichlet $L$-级数之积. 这正是 Artin 当初定义他的映射的原因.

可以证明: 若 $(V,\rho)=\Ind_{H}^{G}(V_{0},\rho_{0})$, 则
\[
L(s,\rho)=L(s,\rho_{0}).
\]
为处理更一般的 Artin $L$-级数, Artin 证明了: 有限群 $G$ 的每个特征都是循环子群的诱导特征的 $\Q$-线性组合. 故
\[
L(s,\rho)=\prod(\text{Dirichlet } L \text{-级数})^{r_{i}},\qquad r_{i}\in\Q.
\]
后来 Brauer 证明了一个更强的定理, 由此可以证明
\[
L(s,\rho)=\prod(\text{Dirichlet } L \text{-级数})^{r_{i}},\qquad r_{i}\in\Z.
\]
Artin 猜想: 只要 $\rho$ 不包含平凡表示, $L(s,\rho)$ 就延拓为整个复平面上的全纯函数. 最后一个公式蕴含: 当诸 $r_i$ 都是正整数时这是对的. 这一猜想在 Langlands 于 $V$ 的维数为 $2$ 的许多情形下证明它之前进展甚微 (见 R.~Langlands, \emph{Base Change for GL(2)}, Princeton, 1980).

\subsection*{Hecke $L$-级数}

Hecke (或 Gr\"ossen) 特征是指从 $\mathbb{I}$ 到 $\CC^{\times}$ 的单位圆的连续同态 $\psi$, 满足 $\psi(K^{\times})=1$, 且对某个有限集 $S$, $\psi$ 在集合 $\{(a_{v})\mid a_{v}=1\ (v\in S),\ a_{v}\ \text{为单位}\ (\text{对所有 } v)\}$ 上取值 $1$.

\begin{example}\label{ex:viii.10.1}
设 $D\in\Z$ 无立方因子, $\zeta$ 为本原三次单位根. 若 $p\equiv 2 \bmod 3$, 则 $p$ 在 $\Q[\zeta]$ 中保持素, 我们令 $\psi(p)=1$. 若 $p\equiv 1 \bmod 3$, 则 $p=\pi\bar{\pi}$ 分解于 $\Q[\zeta]$ 中, 我们取 $\pi$ 满足 $\pi\equiv 1 \bmod 3$. 则 $\pi=\frac{1}{2}(a+3b\sqrt{-3})$, 且 $4p=4\pi\bar{\pi}=a^{2}+27b^{2}$. 现在存在 Hecke 特征使得: 对所有奇数且 $\not\equiv 1 \bmod 3$ 的 $p$, $\psi(p)=1$; 而当 $p\equiv 1\bmod 3$ 时 $\psi(p)=\frac{\pi}{\sqrt{p}}\left(\frac{D}{\pi}\right)$.
\end{example}

对这样的特征, 定义
\[
L(s,\psi)=\prod_{v\notin S}\frac{1}{1-\psi(\pi_{v})\,\mathbb{N}\mathfrak{p}_{v}^{-s}},
\]
其中 $\pi_v$ 是在第 $v$ 个位置为素元、其余位置为 $1$ 的 idèle. Hecke $L$-级数的基本解析性质 (亚纯延拓、函数方程等) 已被充分理解 (例如 J.~Tate, \emph{Fourier analysis in number fields and Hecke's zeta function}, 博士论文, 1950; 转载于 \cite{CasselsFrohlich1967}; 或 \cite{RamakrishnanValenza1999}).

\subsection*{Weil 群与 Artin--Hecke $L$-级数}

关于这一论题, 见 \cite{Tate1979}.

\section*{Gauss 的一个定理}\label{sec:VIII.gauss}

本课程以 Gauss 首先证明的一个定理——二次互反律——开始, 以他的另一个定理结束似乎是恰当的.

考虑椭圆曲线 $E\colon X^{3}+Y^{3}+Z^{3}=0$. 设 $N_p$ 为 $E$ 上坐标在 $\F_p$ 中的点的个数. Gauss 证明了:

(a) 若 $p\not\equiv 1 \bmod 3$, 则 $N_p=p+1$;

(b) 若 $p\equiv 1 \bmod 3$, 则 $N_p=A$, 其中 $A$ 是满足 $4p=A^{2}+27B^{2}$ 且 $A\equiv 1 \bmod 3$ 的唯一整数.

见 \cite{SilvermanTate1992}, IV.2.

Gauss 定理蕴含 Weil 猜想对 $E/\F_p$ 成立, 即
\[
|N_{p}-p-1|<2\sqrt{p}.
\]
它还蕴含 Taniyama 猜想对 $E/\Q$ 成立, 因为它表明 $L$-级数 $L(s,E)$ 等于 $L(s-\frac{1}{2},\psi)$, 其中 $\psi$ 是上例中与 $D=-1$ 相关联的 Hecke 特征.

\begin{quote}
他写信告诉我, 代数数论是他所遇到过的最美的数学主题, 而他苦难中的唯一慰藉就是讲授类域论……这确实是他最钦佩的那类数学: 主要结果范围广阔、美感十足, 但证明在技术上极其艰难.

A.~Borel, 载于: \emph{Current Trends in Mathematics and Physics: A Tribute to Harish-Chandra}, 编辑 S.D.~Adhikari, Narosa Publishing House, New Delhi, Madras, 1995, p.~213.
\end{quote}


% ===== VIII 附录 A/B/C (练习、解答、历史文献来源) =====

\cleardoublepage\phantomsection
\section*{A 附录: 练习}\label{sec:VIII.A}
\addcontentsline{toc}{chapter}{A 附录: 练习}

这些是学员需要提交的练习.

\medskip\noindent\textbf{A-1}\quad
(a) 设 $K$ 是 $\Q$ 的有限 Galois 扩张. 证明 $\Gal(K/\Q)$ 由 $\mathcal{O}_K$ 中素理想的惯性群生成. (提示: 证明由这些惯性群生成的子群的固定域在 $\Z$ 的所有素理想处均不分歧. 素理想 $\mathfrak{P}$ 的惯性群是指分解群 $D(\mathfrak{P})$ 中固定剩余域的那些元素组成的子群; 它的阶等于 $\mathfrak{P}$ 的分歧指数.)

(b) 设 $p_1,\ldots,p_t$ 为互不相同的奇素数, 并令 $p_i^{*}=(-1)^{\frac{p_i-1}{2}}p_i$ (于是只有 $p_i$ 在 $\Q[\sqrt{p_i^{*}}]$ 中分歧). 设 $K=\Q[\sqrt{d}]$, $d=\prod p_i^{*}$.

设 $C_+$ 为 $K$ 的窄类群, 并设 $L$ 为 $C_+^{2}$ 的类域 (于是 $L$ 在 $\mathcal{O}_K$ 的所有素理想处不分歧, 且 $\Gal(L/K)$ 是初等 Abel $2$-群; 此外, $L$ 是使这些陈述成立的最大的域).

证明 $L$ 在 $\Q$ 上是 Galois 的, 并且 $\mathcal{O}_L$ 的每个素理想在 $\Gal(L/\Q)$ 中的惯性群的阶为 $1$ 或 $2$. 由此推出 $\Gal(L/\Q)$ 是初等 Abel $2$-群, 且 $L=\Q[\sqrt{p_1^{*}},\ldots,\sqrt{p_t^{*}}]$. 再由此推出 $(C_+:C_+^{2})=2^{t-1}$ (对照例 \ref{ex:0.7}).

\medskip\noindent\textbf{A-2}\quad
(a) 设 $d>3$ 是满足 $d\equiv 3 \bmod 8$ 的无平方因子整数. 证明存在 $\Q$ 上的三次域 $L$, 其判别式为 $-d$ 或 $-4d$.

(b) 求使 $|\Delta_{L/K}|$ 最小的 $\Q$ 上三次扩张 $L$.

(你可以引用类域论的主要结论.)

\textbf{提示:} 令 $K=\Q[\sqrt{-d}]$. 对于 $3$ 整除 $K$ 的类数的情形, 琢磨课上讲过的 $\Q[\sqrt{-31}]$ 的 Hilbert 类域这个例子. 对于相反的情形, 证明 $2\mathcal{O}_K$ 是一个素理想 $\mathfrak{p}$, 且 $3$ 整除射线类群 $C_{\mathfrak{p}}$ 的阶.

\medskip\noindent\textbf{A-3}\quad 完成练习 \ref{exc:v.3.15} (本书 p.\pageref{exc:v.3.15}). 注意, 在该练习的陈述中, $i$ 被含混地用来同时表示从 $K_{\mathfrak{m},1}$ 到 $I^{S(\mathfrak{m})}$ 与到 $I^{S\cup S(\mathfrak{m})}$ 的 (显然的) 同态.

\medskip\noindent\textbf{A-4}\quad 证明或反驳: 数域 $K$ 的 id\`ele 类群 $\mathbf{C}_K=\mathbb{I}_K/K^{\times}$ 的每个有限指数子群都是开的. (这是练习 \ref{exc:v.5.11} (本书 p.\pageref{exc:v.5.11}). 对 $K_v^{\times}$ 的类似陈述可以引用 —— 见 I \ref{para:1.7}.)

\medskip\noindent\textbf{A-5}\quad 设 $G$ 是一个群, $A$ 是一个带 $G$-作用的群 (不必交换), 即给定一个同态 $G\to\mathrm{Aut}(A)$. 与交换情形一样, 交叉同态 $\varphi\colon G\to A$ 由条件 $\varphi(\sigma\tau)=\varphi(\sigma)\cdot\sigma\varphi(\tau)$ 定义. 若存在 $a\in A$ 使得 $\varphi'(\sigma)=a^{-1}\cdot\varphi(\sigma)\cdot\sigma a$ 对所有 $\sigma\in G$ 成立, 则称两个这样的同态 $\varphi$ 与 $\varphi'$ 等价, 等价类之集记作 $H^{1}(G,A)$ (它是一个带特殊元素之集). 与 Hilbert 定理 90 的证明相类似的论证表明: 当 $L/K$ 是以 $G$ 为 Galois 群的有限 Galois 扩张时, $H^{1}(G,\mathrm{GL}_n(L))=0$.

设 $K$ 是特征 $\ne 2$ 的域. $K$ 上的\textbf{二次空间}是二元组 $(V,\Phi)$, 由 $K$ 上有限维向量空间 $V$ 与 $V$ 上非退化二次型 $\Phi$ 组成, 即 $\Phi$ 是由 $V$ 上非退化对称双线性型 $\varphi$ 伴随的映射 $x\mapsto\varphi(x,x)$. \textbf{等距} $\alpha\colon(V,\Phi)\to(V',\Phi')$ 是向量空间的同构 $\alpha\colon V\to V'$, 满足 $\Phi'(\alpha(v))=\Phi(v)$ 对所有 $v\in V$ 成立.

\textbf{练习:} 设 $L/K$ 是特征 $\ne2$ 的域的有限 Galois 扩张, 并令 $G=\Gal(L/K)$. 取定 $K$ 上的一个二次空间 $(V,\Phi)$, 并令 $O(\Phi)$ 为 $(V,\Phi)\otimes_K L$ 的自同构群 —— 它带有显然的 $G$-作用. 若 $(V',\Phi')$ 是 $K$ 上的二次空间, 且 $\alpha\colon(V,\Phi)\otimes_K L\to(V',\Phi')\otimes_K L$ 是 $L$ 上二次空间的同构, 证明 $\sigma\mapsto\alpha^{-1}\circ\sigma\alpha$ 是交叉同态 $\varphi\colon G\to O(\Phi)$, 其等价类不依赖于 $\alpha$ 的选取. 由此推出 $H^{1}(G,O(\Phi))$ 对在 $L$ 上变得与 $(V,\Phi)$ 同构的 $K$ 上二次空间的同构类之集作出分类 (即 $H^{1}(G,O(\Phi))$ 与该集之间存在一一对应). 再由此推出: 当 $K=\R$, $L=\CC$ 时 (设 $\dim V\ne0$), $H^{1}(G,O(\Phi))\ne0$. (你可以引用 $H^{1}(G,\mathrm{GL}_n(L))=0$.)

\medskip\noindent\textbf{A-6}\quad 设 $H$ 是群 $G$ 中有限指数的子群, 并选取映射 $s\colon H\backslash G\to G$ 使 $H\,s(x)=Hx$, 即为 $G$ 中 $H$ 的每个右陪集取一个代表元 —— 于是 $G=\bigcup_{x\in H\backslash G}H\,s(x)$.

(a) 对 $g\in G$, 证明 $s(x)\cdot g\cdot s(xg)^{-1}\in H$.

对 $g\in G$, 定义
\[
V(g)=\prod_{x\in H\backslash G}s(x)\cdot g\cdot s(xg)^{-1}\ \bmod\ H^{c},
\]
其中 $H^{c}$ 是 $H$ 的换位子群. (注意, 由于 $H/H^{c}$ 交换, $V(g)$ 不依赖于乘积中 $H\backslash G$ 的取序方式.)

(b) 证明 $V(g)$ 不依赖于 $s$ 的选取, 并且它是同态 $G\to H/H^{c}$.

于是 $V$ 诱导同态 $\mathrm{Ver}\colon G/G^{c}\to H/H^{c}$ —— 即 \textbf{Verlagerung} (或\textbf{转移}) 映射.

(c) 对任意 $G$-模 $M$, 证明映射 $m\mapsto N(m)=\sum_{x\in H\backslash G}s(x)\cdot m$ 诱导良定义的同态 $\mathrm{Res}\colon H_0(G,M)\to H_0(H,M)$.

(d) 证明图表
\[
\begin{tikzcd}
G/G^{c} \ar[r, "\simeq"] \ar[d, "\mathrm{Ver}"'] & I_G/I_G^{2} \ar[d, "N"] \\
H/H^{c} \ar[r, "\simeq"'] & I_G/I_H\, I_G \\
H/H^{c} \ar[r, "\simeq"'] & I_H/I_H^{2} \ar[u]
\end{tikzcd}
\]
交换 (水平的同构就是引理 \ref{lem:ii.2.6} 的同构; 向上的映射由包含映射诱导).

[假设 $\mathrm{Res}\colon H_r(G,M)\to H_r(H,M)$ 存在, 与边缘同态相容, 并且在 $r=0$ 时有 (c) 中所给的描述, 则上式表明它们在 $H_1(G,\Z)$ 上与 Verlagerung 一致. 更精确地, 上面的图表可以等同于
\[
\begin{tikzcd}
H_1(G,\Z) \ar[r, "\delta"] \ar[d, "\mathrm{Res}"'] & H_0(G, I_G) \ar[d, "\mathrm{Res}"] \\
H_1(H,\Z) \ar[r, "\delta"] & H_0(H, I_G) \\
H_1(H,\Z) \ar[r, "\delta"'] \ar[u, "\mathrm{id}"'] & H_0(H, I_H) \ar[u]
\end{tikzcd}
\]
边缘同态来自 $G$ 与 $H$ 的增广正合列. 上、下两个边缘同态是同构. 中间的边缘同态是单射, 并且当 $G$ 有限、而 $H_0(H,I_G)$ 换成 $\widehat{H}_T^{-1}(H,I_G)$ 时, 它成为同构.]

\medskip\noindent\textbf{A-7}\quad 本练习将表明上同调与逆向极限不交换 —— 事实上, 映射 $H^{1}(G,\varprojlim M_n)\to\varprojlim H^{1}(G,M_n)$ 不必是单射.

设 $G$ 是一个 profinite 群, 且
\[
\cdots\longrightarrow M_n\overset{\pi_{n-1,n}}{\longrightarrow}M_{n-1}\longrightarrow\cdots\longrightarrow M_1
\]
是一个 $G$-模的投影系统. 假设:

(a) 每个映射 $\pi_{n-1,n}\colon M_n\to M_{n-1}$ 都是满射;

(b) 存在交换图表
\[
\begin{tikzcd}[row sep=small]
\cdots \ar[r] & M_n^{G} \ar[r, "\pi_{n-1,n}"] \ar[d, "\sigma_n"', two heads] & M_{n-1}^{G} \ar[r] \ar[d, "\sigma_{n-1}"', two heads] & \cdots \\
\cdots \ar[r] & \Z \ar[r] & \Z \ar[r] & \cdots
\end{tikzcd}
\]
其中每个竖直映射都是满射.

对任意 $c_n\in\Z$, 级数 $\sum p^{n-1}c_n$ 将在 $\Z_p$ 中收敛 —— 选取序列 $c_n$, $n\ge1$, 使得 $\sum p^{n-1}c_n$ 收敛到 $\Z_p$ 中不属于 $\Z$ 的元素 $c$. 对每个 $n$, 选取元素 $x_{n,1},\ldots,x_{n,n}$ 使得:

(a) 对 $1\le i\le n-1$, $\pi_{n-1,n}(x_{n,i})=x_{n-1,i}$;

(b) $x_{n,n}\in M_n^{G}$ 且 $\sigma_n(x_{n,n})=c_n$.

定义
\[
x_n=x_{n,1}+\cdots+x_{n,n},\qquad
\varphi_n(\sigma)=\sigma x_n-x_n\quad(\sigma\in G).
\]
于是 $\varphi_n$ 是主交叉同态 $G\to M_n$. 注意 $\pi_{n-1,n}(\varphi_n(\sigma))=\varphi_{n-1}(\sigma)$ 对所有 $n$ 成立, 从而存在元素 $\varphi(\sigma)\in M\overset{\mathrm{def}}{=}\varprojlim M_n$, 它在每个 $M_n$ 中的像为 $\varphi_n(\sigma)$. 显然 $\sigma\mapsto\varphi(\sigma)$ 是交叉同态 $G\to M$.

证明 $\varphi$ 不是主交叉同态, 即证明不存在 $M$ 中的元素 $y=(y_n)$ 使得对每个 $n$ 与每个 $\sigma\in G$, 都有 $\varphi_n(\sigma)=\sigma y_n-y_n$.

提示: 令 $y_n=x_n+z_n$, 并证明 $z_1$ 被 $G$ 固定, 但 $\sigma_1(z_1)=c\notin\Z$.

你能构造出满足这些条件的系统 $(M_n,\pi_{n-1,n},\sigma_n)$ 吗?

\medskip\noindent\textbf{A-8}\quad 完成练习 \ref{ex:iv.5.7} 的 (a)--(d) 部分 (若你精力充沛, 也可以完成整个练习). 必要时可以使用第 \ref{ch:IV} 章的非上同调结果.

\medskip\noindent\textbf{A-9}\quad 设 $L_n=\Q_p[\zeta_{p^n}]$, 其中 $\zeta_{p^n}$ 是本原 $p^n$ 次单位根. 证明 $p$ 是来自 $L_n$ 的范数. 由此推出 $L_\infty=\bigcup L_n$ 是 $\Q_p^{\mathrm{ab}}$ 中被 $\mathrm{rec}_{\Q_p}(p)$ 固定的子域 (并在 $\Q_p$ 的某个取定的代数闭包内取并).

\medskip\noindent\phantomsection\hypertarget{item:viii.A-10}{}\textbf{A-10}\quad 设 $L/K$ 是 $n$ 次有限循环数域扩张. 证明: 对每个 $d\mid n$, 满足 $e(\mathfrak{P}/\mathfrak{p})=1$ 且 $f(\mathfrak{P}/\mathfrak{p})=d$ (对某个 —— 因而是所有 —— 位于 $\mathfrak{p}$ 之上的 $\mathfrak{P}$) 的 $K$ 的素 $\mathfrak{p}$ 之集有极密度 $\varphi(d)/n$. (提示: 尝试对 $d$ 作归纳.)

精力充沛者可以类似地给出 Frobenius 密度定理 (Chebotarev 密度定理的较弱、更早的形式) 的一个同样初等的证明: 设 $L/K$ 是以 $G$ 为 Galois 群的 Galois 扩张, 并设 $\sigma$ 是 $G$ 中 $n$ 阶元素. 证明
\[
\{\mathfrak{p}\mid \mathfrak{p}\ \text{不分歧且}\ (\mathfrak{p},L/K)=\sigma^{k}\ \text{对某个与}\ n\ \text{互素的}\ k\}
\]
有极密度 $c\varphi(n)/[L:K]$, 其中 $c$ 是 $\langle\sigma\rangle$ 在 $G$ 中的共轭个数. 再证明 $c\varphi(n)$ 恰是 $\sigma$ 的共轭个数, 从而这与 Chebotarev 密度定理相容. 至于提示, 见 \cite[pp.~207--208]{Marcus1977}.

\medskip\noindent\textbf{A-11}\quad 证明 $16$ 在 $\R$ 中以及在所有奇数 $p$ 的 $\Q_p$ 中是 $8$ 次幂 (但在 $\Q$ 中不是).

[提示: 证明 $\Q[\zeta_8]=\Q[i,\sqrt{2}]$ 在所有奇数 $p$ 处不分歧, 并由此推出: 对奇数 $p$, $\Q_p$ 包含 $1+i$, $\sqrt{2}$, $\sqrt{-2}$ 三者之一.]

注: 这与 Grunwald 定理相矛盾, 而该定理曾被广泛使用了十五年.

\cleardoublepage\phantomsection
\section*{B 附录: 练习解答}\label{sec:VIII.B}
\addcontentsline{toc}{chapter}{B 附录: 练习解答}

\noindent\textbf{练习 A-1 的解答.}

(a) 设 $\mathfrak{p}$ 是 $\mathcal{O}_K$ 中的素理想, 并设 $I(\mathfrak{p})$ 是它的惯性群. 则 $M=K^{I(\mathfrak{p})}$ 具有如下性质: $\mathfrak{p}\cap\mathcal{O}_M$ 在 $\mathfrak{p}\cap\Z$ 上不分歧 (ANT, 8.11). 于是被所有惯性群固定的任何域都在 $\Q$ 上不分歧, 从而等于 $\Q$ (ANT, 4.9).

(b) 域 $L$ 在 $\Q$ 上是 Galois 的, 因为它在它的 Galois 闭包的任何自同构下稳定.

显然惯性群的阶为 $1$ 或 $2$, 于是 $\Gal(L/\Q)$ 由 $2$ 阶元素生成. 这并不意味着它是交换的 (每个二面体群都由两个 $2$ 阶元素生成), 但存在扩张
\[
0\to\Gal(L/K)\to\Gal(L/\Q)\to\Gal(K/\Q)\to 0,
\]
其中 $\Gal(L/K)\simeq(\Z/2\Z)^{s}$ ($s$ 为某个数), $\Gal(K/\Q)\simeq\Z/2\Z$, 而非平凡惯性群的生成元在 $\Gal(K/\Q)$ 中的像 $\ne 1$ (否则 (a) 中的论证将表明分歧发生在 $L/K$ 而非 $K/\Q$ 中). 设 $\sigma$ 与 $\tau$ 生成两个 $2$ 阶惯性群. 则 $\sigma|_K\ne1$ 且 $\tau|_K\ne1$, 于是 $\sigma\tau|_K=1$. 从而 $\sigma\tau\in\Gal(L/K)$, 故其阶为 $1$ 或 $2$. 因此 $\sigma$ 与 $\tau$ 交换, 于是 $\Gal(L/\Q)$ 由互相交换的 $2$ 阶元素生成, 这表明它是初等 Abel $2$-群.

在 $L$ 中只有素 $p_1,\ldots,p_t$ 分歧, 且对应于每个 $p_i$ 只有一个惯性群 (在 $p_i$ 之上可能有许多 $L$ 的素理想, 但由于 $\Gal(L/\Q)$ 交换, 每个的惯性群都相同). 于是 $\Gal(L/\Q)\simeq(\Z/2\Z)^{s}$, $s\le t$. 由于 $L$ 包含 $\Q[\sqrt{p_1^{*}},\ldots]$, 而后者在 $\Q$ 上的次数为 $2^{t}$, 可见 $L=\Q[\sqrt{p_1^{*}},\ldots]$. 其余是容易的.

[关于 $L=\Q[\sqrt{p_1^{*}},\ldots]$ 的另一种证明: 由于 $\Gal(L/\Q)$ 是初等 Abel $2$-群, $L=\Q[\sqrt{m_1},\sqrt{m_2},\ldots]$, 其中 $m_i\in\Z$ (容易的 Kummer 理论); 但容易看出, $\Q[\sqrt{p_1^{*}},\ldots]$ 是这类域中其与 $K$ 的复合在 $K$ 上不分歧的最大者.]

\noindent\textbf{练习 A-2 的解答.}

由于 $d\equiv 3 \bmod 8$, 有 $-d\equiv 1 \bmod 4$, 从而 $\Delta_{K/\Q}=-d$. 我用``素''一词指``有限素''.

若 $3$ 整除 $K$ 的类数, 则存在 $K$ 的 $3$ 次不分歧扩张 $L$ (由类域论). 与 $d=31$ 的情形一样, $L$ 在 $\Q$ 上是 Galois 的, 且 $L=M\cdot K$, 其中 $M$ 是 $\Q$ 上 $3$ 次且非 Galois 的域. 由 Galois 理论, 含于 $M$ 的 Galois 闭包中的 $\Q$ 的唯一二次扩张是 $\Q[\sqrt{\Delta_{M/\Q}}]$, 于是 $\Delta_{M/\Q}=-d\times$平方 (容易的 Kummer 理论). 若素数 $p$ 不整除 $d$, 则 $p$ 在 $L$ (从而在 $M$) 中不分歧, 于是 $p$ 不整除 $\Delta_{M/K}$. 若 $p\mid d$, 则 $p\mathcal{O}_K=\mathfrak{p}_1^{2}\mathfrak{p}_2$, 这 (差理论) 表明 $p$ 整除 $\Delta_{M/K}$ 但不整除 $p^{2}$. 因此 $\Delta_{M/\Q}=-d$.

假设 $3$ 不整除类数. 由于 $-d\equiv 5 \bmod 8$, $(2)$ 在 $\mathcal{O}_K$ 中保持素\footnote{若 $-d\equiv 1 \bmod 8$, 它就会分裂.}, 记 $2\mathcal{O}_K=\mathfrak{p}$, 且 $3$ 整除射线类群 $C_{\mathfrak{p}}$ 的阶 (射线类群包含 $(\mathcal{O}_K/\mathfrak{p})^{\times}=\mathbb{F}_4^{\times}$ —— 见 V \ref{thm:v.1.7}; 这里我们利用 $d>3$ 来看出 $K$ 的单位只有 $\pm1$). 我们用与前面相同的论证, 取 $L$ 为 $K$ 的仅在 $\mathfrak{p}$ 处分歧的 $3$ 次扩张. 同样 $L=M\cdot K$, 其中 $M$ 与前面一样, 只是 $2$ 在 $M$ 中全分歧, 论证表明 $\Delta_{M/\Q}=-4d$.

从本书 p.\pageref{exc:v.3.14} 的表可见 $\Delta_{M/\Q}=-23$ 是可能的. 由 Minkowski 界, $|\Delta_{M/\Q}|\ge 13$ (有复素) 或 $|\Delta_{M/\Q}|\ge 21$ (无复素). Stickelberger 定理排除了其余情形, 只留下 $-15,-16,-19,-20,21$. 若 $M$ 的判别式 $d=-15,-19,21$, 则其 Galois 闭包 $L=M\cdot\Q[\sqrt{d}]$ 在 $\Q[\sqrt{d}]$ 上不分歧 (否则 $\Q$ 的某个素会在 $L$ 中全分歧而在 $M$ 中不然), 但 $\Q[\sqrt{d}]$ 的类数分别为 $2,1,1$. 对 $d=-16,-20$, 同样的论证表明只有位于 $2$ 之上的素可能在 $L=M\cdot\Q[\sqrt{d}]$ 中分歧, 但此处类数为 $1$ 或 $2$, 且相关的射线类数不被 $3$ 整除 ($2$ 在 $\Q[\sqrt{-1}]$ 与 $\Q[\sqrt{-5}]$ 中分歧).

\textbf{注:}

(a) 可以不用差理论, 而使用 ANT 3.39(a) 中的公式 (但该公式的证明用到了差理论).

(b) 用手计算上述类数并不太难, 但简便的办法是使用 PARI.

\noindent\textbf{练习 A-3 的解答.}

在本讲义 (及本解答) 中, $i$ 被含混地用来同时表示从 $K_{\mathfrak{m},1}$ 到 $I^{S(\mathfrak{m})}$ 与到 $I^{S\cup S(\mathfrak{m})}$ 的 (显然的) 同态. 这些映射使图表
\[
\begin{tikzcd}
 & & K_{\mathfrak{m},1} \ar[d] \ar[drr] & & \\
\Z S \ar[rr] & & I^{S(\mathfrak{m})} \ar[rr] & & I^{S\cup S(\mathfrak{m})}
\end{tikzcd}
\]
交换 —— 其中 $\Z S$ 是由 $S$ 生成的自由 Abel 群. 该三角形的核--余核正合列给出
\[
\Z S\to I^{S(\mathfrak{m})}/i(K_{\mathfrak{m},1})\to I^{S\cup S(\mathfrak{m})}/i(K_{\mathfrak{m},1})\to 0,
\]
于是
\[
I^{S(\mathfrak{m})}/\langle S\rangle\cdot i(K_{\mathfrak{m},1})\simeq I^{S\cup S(\mathfrak{m})}/i(K_{\mathfrak{m},1}).
\]
因此, $I^{S\cup S(\mathfrak{m})}$ 中包含 $i(K_{\mathfrak{m},1})$ 的子群与 $I^{S(\mathfrak{m})}$ 中包含 $\langle S\rangle\cdot i(K_{\mathfrak{m},1})$ 的子群在对应 $H\leftrightarrow H'$ 之下自然一一对应. 对这样的一对 $H$, $H'$, 设 $L$ 是 $H'$ 的类域. 我们得到同构
\[
I^{S\cup S(\mathfrak{m})}/H\to I^{S(\mathfrak{m})}/H'\to\Gal(L/K).
\]
第二个映射是 Artin 映射, 由 $\mathfrak{p}\mapsto(\mathfrak{p},L/K)$ 诱导. 当然 $L/K$ 是 Abel 的, 且 $S$ 中的素分裂, 因为它们属于 $H'$ (从而有 $(\mathfrak{p},L/K)=1$). 此外, 不在 $S(\mathfrak{m})$ 中的素不分歧. 最后, 两个映射的复合也是 Artin 映射, 于是不在 $S\cup S(\mathfrak{m})$ 中且分裂的素恰好是 $H$ 中的那些. 这表明 $L/K$ 是 $H$ 的 $S$-类域. 我们已经证明它有 Galois 群 $I^{S\cup S(\mathfrak{m})}/H$, 并且它是唯一的, 因为我们知道哪些素在其中分裂 (除有限个可能的例外).

反之, 任何满足 (a) 的 $L$ 都是某个使 $S(\mathfrak{m})\cap S=\emptyset$ 的模 $\mathfrak{m}$ 以及 $I^{S(\mathfrak{m})}$ 的包含 $\langle S\rangle\cdot i(K_{\mathfrak{m},1})$ 的某个子群 $H'$ 的类域, 由此可知, 它是 $I^{S\cup S(\mathfrak{m})}$ 中与 $H'$ 对应的子群 $H$ 的 $S$-类域.

\noindent\textbf{练习 A-4 的解答.}

唉, id\`ele 类群的每个有限指数子群都是开的这一命题并不成立, 即使对数域也如此. 如同 I \ref{cor:1.5} 中那样, 设 $V$ 是可数无穷个 $\mathbb{F}_p$ 拷贝的乘积 (或任何其他有限域的拷贝的乘积), 并设 $W\subset V$ 由那些除有限多个分量外全为零的元素组成 (于是 $W$ 是 $\mathbb{F}_p$ 拷贝的直和). 赋予乘积拓扑, $V$ 是紧的 (且不可数). 子空间 $W$ 可数, 但在 $V$ 中稠密 —— 事实上, 它满射到 $V$ 的任何有限商. 商 $V/W$ 是 $\mathbb{F}_p$ 上的向量空间, 因而有基 $B$ (这需要选择公理 —— 见 Jacobson, \emph{Lectures in Abstract Algebra}, Vol II, Chap. IX; 说 $B$ 是基, 意指 $B$ 的每个有限子集都线性无关, 且 $V/W$ 的每个元素都是 $B$ 中元素的有限线性组合). 设 $S$ 是 $B$ 的非空有限子集, 并设 $W'$ 是 $\mathrm{span}(B\setminus S)$ 在 $V$ 中的原像. 则 $V\supset W'\supset W$, 且 $W'$ 是 $V$ 中指数有限的真子群, 其闭包为 $V$. 它不可能是开的, 因为否则它将是闭的.

回忆 (本讲义 \ref{lem:v.5.9}) $\mathbf{C}_{\Q}\simeq\R_{>0}\times\prod\Z_p^{\times}$ (拓扑群同构). 每个 $\Z_p^{\times}$ 有指数 $2$ 的开子群. 取相应商映射的乘积, 我们得到连续同态 $\prod_p\Z_p^{\times}\to V$, 其中 $V$ 是以素数为指标集的 $\mathbb{F}_2$ 拷贝的乘积. 设 $W'$ 是 $V$ 中不闭 (从而不开) 的有限指数子群. 它在 $\prod\Z_p^{\times}$ 中的原像仍不闭且指数有限, 它在 $\mathbf{C}_{\Q}$ 中的原像亦然.

类似的论证表明 $\Gal(\Q^{\mathrm{al}}/\Q)$ 有不开的有限指数子群 —— 考虑其商群 $\Gal(K/\Q)$, 其中 $K=\Q[\sqrt{2},\sqrt{3},\sqrt{5},\ldots]$ (FT 7.26).

注意, 在每种情形我们都用了选择公理来构造不开的有限指数子群. 能否避免选择公理并找到显式的例子?

\noindent\textbf{练习 A-5 的解答.}

对在左 $G$-作用之下的对象之间的映射 $\alpha$, 定义 $\sigma\alpha\overset{\mathrm{def}}{=}\sigma\circ\alpha\circ\sigma^{-1}$ (这是标准的). 显然 $\sigma(\alpha\circ\beta)=\sigma\alpha\circ\sigma\beta$.

设 $V$ 是 $K$ 上的向量空间. 若线性映射 $\alpha\colon V\otimes L\to V\otimes L$ 相对于 $V$ 的某组基有矩阵 $(a_{ij})$, 则 $\sigma\alpha$ 有矩阵 $(\sigma a_{ij})$. 因此, $\alpha$ ``在 $K$ 上定义'', 即形如 $\alpha_0\otimes1$, 当且仅当 $\sigma\alpha=\alpha$ 对所有 $\sigma\in G$ 成立. 对二次型也有类似的结论.

在 $O(\Phi)$ 中, $\cdot$ 定义为 $\circ$.

现在设 $\alpha\colon(V,\Phi)\otimes L\to(V',\Phi')\otimes L$ 是二次空间的同构, 如练习中所述. 令 $\varphi(\sigma)=\alpha^{-1}\circ\sigma\alpha$. 则
\[
\varphi(\sigma)\cdot\sigma\varphi(\tau)=\alpha^{-1}\circ\sigma\alpha\circ\sigma(\alpha^{-1}\circ\tau\alpha)=\alpha^{-1}\circ\sigma\alpha\circ\sigma\alpha^{-1}\circ\sigma\tau\alpha=\varphi(\sigma\tau).
\]
因此 $\varphi$ 是交叉同态. 任何其他同构 $\alpha'\colon(V,\Phi)\otimes L\to(V',\Phi')\otimes L$ 与 $\alpha$ 相差 $(V,\Phi)\otimes L$ 的一个自同构: $\alpha'=\alpha\circ\beta$, $\beta\in O(\Phi)$. 由 $\alpha'$ 定义的交叉同态是 $\varphi'$, 满足
\[
\varphi'(\sigma)=\beta^{-1}\circ\alpha^{-1}\circ\sigma(\alpha\circ\beta)=\beta^{-1}\circ\varphi(\sigma)\circ\sigma\beta,
\]
这表明 $\varphi'$ 与 $\varphi$ 等价.

于是每个 $(V',\Phi')$ 定义了 $H^{1}(G,O(\Phi))$ 的一个元素 $[(V',\Phi')]$, 且不难验证: 若 $(V',\Phi')\simeq(V'',\Phi'')$, 则 $[(V',\Phi')]=[(V'',\Phi'')]$. 反之, 若 $[(V',\Phi')]=[(V'',\Phi'')]$, 则可以选取同构 $\alpha\colon(V,\Phi)\otimes L\to(V',\Phi')\otimes L$ 与 $\beta\colon(V,\Phi)\otimes L\to(V'',\Phi'')\otimes L$ 使它们给出同一个交叉同态 —— 因而 $\alpha^{-1}\circ\sigma\alpha=\beta^{-1}\circ\sigma\beta$ 对所有 $\sigma\in G$ 成立. 此时 $\beta\circ\alpha^{-1}\colon(V',\Phi')\to(V'',\Phi'')$ 是满足 $\sigma(\beta\circ\alpha^{-1})=\beta\circ\alpha^{-1}$ 对所有 $\sigma$ 成立的同构, 从而在 $K$ 上定义.

剩下要证明 $H^{1}(G,O(\Phi))$ 的每个元素都来自某个 $(V',\Phi')$. 设 $\varphi$ 是交叉同态 $G\to O(\Phi)$. 由于 $H^{1}(G,\mathrm{GL}_n)=0$, 有 $\varphi(\sigma)=\alpha^{-1}\circ\sigma\alpha$, 其中 $\alpha$ 是 $V\otimes L$ 的某个自同构. 设 $\Phi'$ 是 $V\otimes L$ 上满足 $\Phi'\circ\alpha=\Phi$ 的二次型. 我断言 $\Phi'$ 在 $K$ 上定义: 事实上,
\[
\sigma\Phi'=\sigma(\Phi\circ\alpha^{-1})=\Phi\circ\sigma\alpha^{-1}=\Phi\circ\varphi(\sigma)^{-1}\circ\alpha^{-1}=\Phi\circ\alpha^{-1}=\Phi'.
\]
二次型 $\Phi'$ 的选取使得 $\alpha$ 是同构 $(V,\Phi)\to(V,\Phi')$, 由此可知 $(V,\Phi')$ 由 $\varphi$ 表示.

对最后的陈述, 注意二次空间 $(\R^{2},x^{2}+y^{2})$ 不同构于二次空间 $(\R^{2},x^{2}-y^{2})$ (前者只取正值), 但在 $\CC$ 上变得与它同构.

\noindent\textbf{练习 A-6 的解答.}

(a) 按定义, $x=H\cdot s(x)$ 且 $xg=H\cdot s(xg)$. 把第一个等式两边乘以 $g$, 得 $xg=H\cdot s(x)g$. 于是 $H\cdot s(x)g=H\cdot s(xg)$, 从而 $s(x)\cdot g\cdot s(xg)^{-1}\in H$.

(b) 设 $(s^{0}(x))$ 是第二组右陪集代表元, 即 $s^{0}(x)=h(x)s(x)$. 则
\[
\prod_x s^{0}(x)\cdot g\cdot s^{0}(xg)^{-1}
=\prod_x h(x)\cdot s(x)\cdot g\cdot s(xg)^{-1}\cdot h(xg)^{-1}
\equiv\prod_x s(x)\cdot g\cdot s(xg)^{-1}\cdot\prod_x h(x)\cdot\prod_x h(xg)^{-1}\ \bmod\ H^{c}.
\]
由于当 $x$ 取遍 $H$ 在 $G$ 中的右陪集时, $xg$ 亦然, 故 $\prod_x h(x)=\prod_x h(xg)$.

对 $g,g^{0}\in G$,
\begin{multline*}
\prod_x s(x)\cdot gg^{0}\cdot s(xgg^{0})^{-1}
\equiv\prod_x s(x)\cdot g\cdot s(xg)^{-1}\cdot\prod_x s(xg)\cdot g^{0}\cdot s(xgg^{0})^{-1}\ \bmod\ H^{c}\\
=\prod_x s(x)\cdot g\cdot s(xg)^{-1}\cdot\prod_x s(x)\cdot g^{0}\cdot s(xg^{0})^{-1},
\end{multline*}
—— 我们再次利用了当 $x$ 取遍 $H$ 的右陪集时, $xg$ 亦然这一事实. 按 $\bmod\ H^{c}$ 读出, 此式成为
\[
V(gg^{0})=V(g)\cdot V(g^{0}).
\]

(c) 设 $(s^{0}(x))_x$ 是第二组右陪集代表元, 即 $s^{0}(x)=h(x)s(x)$. 对 $m\in M$,
\[
\sum_x s^{0}(x)\cdot m=\sum_x s(x)\cdot m+\sum_x(h(x)-1)\cdot s(x)\cdot m\equiv\sum_x s(x)\cdot m\ \bmod\ I_H M.
\]
因此,
\[
m\mapsto\sum_x s(x)\cdot m\colon M\to M/I_H M\qquad(*)
\]
不依赖于 $s$ 的选取.

若 $(s(x))_x$ 是右陪集代表元系, 则对任意固定的 $g\in G$, $x\mapsto s(xg^{-1})g$ 亦然, 于是
\[
\sum_x s(x)\cdot g\, m\equiv\sum_x s(xg)\, m\ \bmod\ I_H M.
\]
因此 $\sum_x s(x)\cdot(g-1)m\equiv 0 \bmod I_H M$, 这表明映射 $(*)$ 经过 $M/I_G M$ 分解.

它显然是交换群的同态.

(d) 如 \cite[p.~129]{Serre1962} 所说: ``Il n'y a plus maintenant qu'\`a calculer''.

设 $g\in G$. 它走捷径在 $I_G/I_H\, I_G$ 中的像是
\[
\sum_x s(x)(g-1)\ \bmod\ I_H\, I_G,
\]
而走远路的像是
\[
\sum_x(s(x)\cdot g\cdot s(xg)^{-1}-1)\ \bmod\ I_H\, I_G.
\]
对每个 $x$, 存在 $h_x\in H$ 使
\[
s(x)\cdot g=h_x\cdot s(xg)
\]
(见解答的第一段). 现在,
\begin{multline*}
\sum_x s(x)(g-1)=\sum_x h_x\cdot s(xg)-\sum_x s(xg)=\sum_x(h_x-1)s(xg)\\
\equiv\sum_x(h_x-1)\ \bmod\ I_H\, I_G=\sum_x(s(x)\cdot g\cdot s(xg)^{-1}-1).
\end{multline*}
或者, 注意到\footnote{更妙的做法 (Darij Grinberg): 注意 $\sum_x s(x)(g-1)=\sum_x s(x)g-\sum_x s(x)=\sum_x s(x)g-\sum_x s(xg)$; 且 $(s(x)\cdot g\cdot s(xg)^{-1}-1)-s(x)\cdot g+s(xg)=(s(x)\cdot g\cdot s(xg)^{-1}-1)(1-s(xg))$.}
\[
\sum_x s(x)(g-1)=\sum_x s(x)g-\sum_x s(x)=\sum_x s(x)g-\sum_x s(xg),
\]
以及
\[
s(x)\cdot g\cdot s(xg)^{-1}-1-s(x)\cdot g+s(xg)=(s(x)\cdot g\cdot s(xg)^{-1}-1)(1-s(xg)),
\]
后者显然属于 $I_H\cdot I_G$.

\noindent\textbf{练习 A-7 的解答.}

假设存在这样的 $y=(y_n)$. 按提示, 写 $y_n=x_n+z_n$. 由于 $\varphi_n=\sigma x_n-x_n$, 故 $\sigma z_n=z_n$ 对所有 $\sigma$ 与 $n$ 成立. 由假设, $\pi_{n-1,n}$ 把 $y_n$ 映到 $y_{n-1}$, 而经计算, 它把 $x_n$ 映到 $x_{n-1}+\pi_{n-1,n}(x_{n,n})$. 因此, 把各个 $\pi$ 逐次应用到定义 $y_n$ 的等式上, 我们得到 (对 $\pi$ 采用显然的记号)
\[
\begin{aligned}
y_{n-1}&=x_{n-1}+\pi_{n-1,n}(x_{n,n})+\pi_{n-1,n}(z_n)\\
y_{n-2}&=x_{n-2}+\pi_{n-2,n}(x_{n,n})+\pi_{n-2,n-1}(x_{n-1,n-1})+\pi_{n-2,n}(z_n)\\
&\;\;\vdots\\
y_1&=x_1+\pi_{1,n}(x_{n,n})+\cdots+\pi_{1,2}(x_{2,2})+\pi_{1,n}(z_n).
\end{aligned}
\]
把它与定义 $z_1$ 的等式比较, 我们发现\footnote{Darij Grinberg 建议在下面两个等式中删去 $x_{1,1}$ 与 $c_1$.}
\[
z_1=x_{1,1}+\pi_{1,2}(x_{2,2})+\cdots+\pi_{1,n}(x_{n,n})+\pi_{1,n}(z_n).
\]
对它应用 $\sigma_1$, 得
\[
\sigma_1(z_1)=c_1+pc_2+\cdots+p^{n-1}c_n+\sigma_1(\pi_{1,n}(z_n)).
\]
现在, $\sigma_1(\pi_{1,n}(z_n))=p^{n-1}\sigma_n(z_n)$, 于是令 $n\to\infty$, 我们发现 $\sigma_1(z_1)=c\notin\Z$, 矛盾.

我们可以如下构造满足条件的系统 $(M_n,\pi_{n-1,n},\sigma_n)$. 设 $G$ 是一个有商 $G_n\to G_{n-1}$ 的 profinite 群, 其中 $G_n$ 的阶为 $p^{n}$ 且 $G_n\to G_{n-1}$ 为满射 (例如 $G=\Z_p$, $G_n=\Z/p^n\Z$), 并令 $M_n=\Z[G_n]$. 映射 $\pi_{n-1,n}$ 由 $G_n\to G_{n-1}$ 诱导. 由于 $M_n^{G}=\Z\cdot(\sum_{\sigma\in G_n}\sigma)$, 我们可以把 $\sigma_n$ 定义为把 $\sum_{\sigma\in G_n}\sigma$ 映到 $1$ 的映射.

\noindent\textbf{练习 A-8 的解答.}

(a) 设 $A$ 是 $F$ 上的中心单代数. 若 $A$ 是矩阵代数, 则 $A\simeq H(1,1)$ (见下面的 (c)). 否则 (Wedderburn 定理) 它是可除代数, 我们现在假设这一点. 设 $i\in A$, $i\notin F$. 则 $F[i]$ 是 $F$ 的二次域扩张. 配方之后, 可设 $i^{2}=a\in F$. 根据 Noether--Skolem 定理, 存在 $j\in A$ 使 $jij^{-1}=\bar{i}$. 现在 $j^{2}$ 与 $F[i]$ 交换, 故 $j^{2}\in F[i]$. 由于 $j^{2}$ 又与 $j$ 交换, 它属于 $F$, 记 $j^{2}=b\in F$. 剩下要验证 $1,i,j,ij$ 线性无关. 设
\[
c+di+ej+fij=0,\qquad c,d,e,f\in F.
\]
用 $i$ 共轭, 得
\[
c+di-ej-fij=0,
\]
于是 $c+di=0$, 从而 $c=0=d$. 再由 $0=ej+fij=(e+fi)j$ 得 $e=0=f$. 这就证明了 $A\simeq H(a,b)$.

(b) 设 $A=H(a,b)$. 若对某组不全为零的 $w,x,y,z\in F$ 有 $w^{2}-ax^{2}-by^{2}+abz^{2}=0$, 则 $\alpha\overset{\mathrm{def}}{=}w+xi+yj+zk$ 是非零零因子, 从而 $A$ 不是可除代数. 反之, 假设对所有非零 $\alpha\in A$ 都有 $\alpha\bar{\alpha}\ne0$. 由于 $\alpha\bar{\alpha}\in F\setminus\{0\}$, 它有逆, 从而 $\alpha$ 可逆, 其逆为 $\bar{\alpha}\cdot(\alpha\bar{\alpha})^{-1}$.

(c) 设 $\alpha=\begin{pmatrix}0&1\\1&0\end{pmatrix}$ 与 $\beta=\begin{pmatrix}1&0\\0&-1\end{pmatrix}$. 直接验证可知 $I,\alpha,\beta,\alpha\beta=\begin{pmatrix}0&-1\\1&0\end{pmatrix}$ 构成 $M_2(F)$ 的一组基, $\alpha^{2}=I$, $\beta^{2}=I$, 且 $\alpha\beta=-\beta\alpha$. 因此 $M_2(F)\simeq H(1,1)$. 或者, 这也可以由 (b) 推出 (利用 Wedderburn 定理).

(d) 映射 $i\mapsto xi$, $j\mapsto yj$ 定义了 $F$-代数同构 $H(a,b)\to H(ax^{2},by^{2})$. (DG 认为映射应为 $i\mapsto i/x$, $j\mapsto j/y$.)

(e) 与 $L$ 作张量积不改变 $H(a,b)$ 的一组基的乘法表.

(f) 设 $A=H(a,b)$. 设 $I$ 是 $A$ 中的真双侧理想. 由于 $F\cap I=\{0\}$, 陪集 $1+I$, $i+I$, $j+I$, $ij+I$ 具有与原元素相同的乘法表, 于是 $A/I$ 是四元数代数 $\simeq H(a,b)$. 特别地, 它在 $F$ 上的次数为 $4$, 这蕴含 $I$ 为零. 因此 $A$ 是单代数. 容易看出 $F$ 是 $A$ 的中心.

(g) 二次型
\[
w^{2}-ax^{2}-(1-a)y^{2}+a(1-a)z^{2}
\]
有非平凡零点 $(t,t,t,0)$, 其中 $t$ 为 $F$ 中任意非零元. 于是可应用 (b).

(h) 二次型
\[
w^{2}-x^{2}-by^{2}+bz^{2}
\]
有非平凡零点 $(t,t,s,s)$, 其中 $s,t\in F$ 不全为零. 二次型
\[
w^{2}-ax^{2}+ay^{2}-a^{2}z^{2}
\]
有非平凡零点 $(at,s,s,t)$, 其中 $s,t\in F$ 不全为零. 于是可应用 (b). (当然, 不假设 Wedderburn 定理而直接证明 (g) 与 (h) 是更有意思的.)

(i) 映射 $1\mapsto1$, $i\mapsto i$, $j\mapsto j$, $k\mapsto-k$ 是同构 $H(a,b)\to H(a,b)^{\mathrm{opp}}$.

(j) 若 $b$ 是 $F$ 中的平方数, 则由 (d) 与 (h) 得 $H(a,b)\simeq M_2(F)$, 且 $a\in\mathrm{Nm}(F[\sqrt{b}])$. 于是假设它不是平方数. 设 $(w,x,y,z)$ 是 $w^{2}-ax^{2}-by^{2}+abz^{2}$ 的非平凡零点. 若 $x^{2}-bz^{2}=0$, 则 $w^{2}-by^{2}=0$, 从而 $w,x,y,z$ 全为零, 矛盾. 因此 $x^{2}-bz^{2}\ne0$, 且
\[
a=\frac{w^{2}-by^{2}}{x^{2}-bz^{2}},
\]
它是来自 $F[\sqrt{b}]$ 的范数 (作为两个范数之商). 反之, 若 $a=w^{2}-by^{2}$, 则 $w^{2}-a\cdot1^{2}-by^{2}+ab\cdot0^{2}=0$, 于是 $(w,1,y,0)$ 是该二次型的非平凡零点.

(k) 用引理 \ref{lem:iv.3.15} 证明 $H(a,b)\otimes_k H(a,c)\approx H(a,bc)$.

\noindent\textbf{练习 A-9 的解答.}

$p^{n}$ 次分圆多项式为
\[
\Phi(X)=\frac{X^{p^{n}}-1}{X^{p^{n-1}}-1}=X^{p^{n-1}(p-1)}+X^{p^{n-1}(p-2)}+\cdots+X^{p^{n-1}}+1.
\]
它的根是本原 $p^{n}$ 次单位根.

令 $F(X)=\Phi(X+1)$. 则 $F(0)=\Phi(1)=p$, 且
\[
F(X)=\frac{(X+1)^{p^{n}}-1}{(X+1)^{p^{n-1}}-1}\equiv\frac{X^{p^{n}}}{X^{p^{n-1}}}=X^{p^{n-1}(p-1)} \bmod p.
\]
因此 $F(X)$ 是 Eisenstein 多项式, 从而在 $\Q_p$ 上不可约. 它的分裂域是 $L_n$. 显然 $L_n=\Q_p[\zeta_{p^n}-1]$, 且 $F(X)$ 是 $\zeta_{p^n}-1$ 的极小多项式. 因此
\[
\mathrm{Nm}_{L_n/\Q_p}(\zeta_{p^n}-1)=(-1)^{\deg F}F(0)=p
\]
(这里利用了 $\deg F=p^{n-1}(p-1)$ 为偶数 —— 除平凡情形 $n=1$, $p=2$ 外均如此).

第二个陈述是显然的.

\noindent\textbf{练习 A-10 的解答.}

我们必须证明: 对每个 $d\mid n$, 在 $L$ 中不分歧且 $(\mathfrak{p},L/K)$ 的阶为 $d$ 的 $K$ 的素 $\mathfrak{p}$ 之集有极密度 $\varphi(d)/n$. 我们对 $d$ 作归纳.

若 $d=1$, 该集由在 $L$ 中完全分裂的素组成, 而我们知道这样的素之集的密度为 $1/n=\varphi(1)/n$.

设 $H$ 是 $G$ 中阶整除 $d$ 的元素组成的子群 —— 它是 $d$ 阶循环子群. 元素 $\sigma\in G=\Gal(L/K)$ 的阶整除 $d$ 当且仅当 $\sigma|_{L^{H}}=\mathrm{id}$. 因此, $(\mathfrak{p},L/K)$ 的阶整除 $d$ 当且仅当 $(\mathfrak{p},L/K)|_{L^{H}}=1$. 但 $(\mathfrak{p},L/K)|_{L^{H}}=(\mathfrak{p},L^{H}/K)$, 于是 $(\mathfrak{p},L/K)$ 的阶整除 $d$ 当且仅当 $\mathfrak{p}$ 在 $L^{H}$ 中分裂 —— 这样的素之集的密度为 $d/n$. 由归纳假设, 其中 $(\mathfrak{p},L/K)$ 的阶为 $d'\mid d$, $d'\ne d$ 的那些素具有密度 $\varphi(d')/n$. 于是, 使 $(\mathfrak{p},L/K)$ 的阶恰好为 $d$ 的素 $\mathfrak{p}$ 之集的密度为
\[
\frac{d}{n}-\sum_{\substack{d'\mid d\\ d'\ne d}}\frac{\varphi(d')}{n}
=\frac{d-\sum\varphi(d')}{n}.
\]
但 $\sum_{d'\mid d}\varphi(d')=d$ (若 $\zeta$ 是本原 $d$ 次单位根, 则对 $d$ 的每个因子 $d'$, 恰有一个幂 $\zeta^{m}$ 是本原 $d'$ 次单位根), 这就完成了证明.

\noindent\textbf{练习 A-11 的解答.}

注意 $\zeta_8=\frac{1+i}{\sqrt{2}}$, 于是 $\Q[\zeta_8]\subset\Q[i,\sqrt{2}]$. 由于二者在 $\Q$ 上的次数相同, 必然相等.

$X^{2}-1$, $X^{2}+2$, $X^{2}-2$ 的判别式都只被 $2$ 整除.

设 $p$ 为奇素数. $-1,2,-2$ 中至少有一个是 $\mathbb{F}_p$ 中的平方数 —— 例如, 若 $-1$ 与 $2$ 都不是平方数, 则它们的乘积 $-2$ 必是 (平方数构成 $\mathbb{F}_p^{\times}$ 中指数为 $2$ 的子群). 因此这些多项式中至少有一个在 $\mathbb{F}_p[X]$ 中分裂, 且由 Hensel 引理也在 $\Q_p[X]$ 中分裂. 这表明 $-1,2,-2$ 中至少有一个是 $\Q_p$ 中的平方数. 由于 $16=(1+i)^{8}=(\sqrt{2})^{8}=(\sqrt{-2})^{8}$, 可知 $16$ 是 $\Q_p$ 中的 $8$ 次幂.

它也是 $\R$ 中的 $8$ 次幂.

\cleardoublepage\phantomsection
\section*{C 附录: 类域论历史的文献来源}\label{sec:VIII.C}
\addcontentsline{toc}{chapter}{C 附录: 类域论历史的文献来源}

\medskip\noindent Artin, Emil; Nesbitt, Cecil J.; Thrall, Robert M.: \emph{Rings with Minimum Condition}. University of Michigan Publications in Mathematics, no.~1. University of Michigan Press, Ann Arbor, Mich., 1944. x+123~pp. MR0010543.\par\medskip

\noindent Artin, Emil; Tate, John: \emph{Class field theory}. AMS Chelsea Publishing, Providence, RI, 2009. viii+194~pp.\par\medskip

\noindent Brauer, Richard: Emil Artin. \emph{Bull. Amer. Math. Soc.} \textbf{73} (1967), 27--43.\par\medskip

\noindent Conrad, Keith: History of class field theory (见作者网页上的在线讲义).\par\medskip

\noindent Fenster, Della D.; Schwermer, Joachim: A delicate collaboration: Adrian Albert and Helmut Hasse and the principal theorem in division algebras in the early 1930's. \emph{Arch. Hist. Exact Sci.} \textbf{59} (2005), no.~4, 349--379.\par\medskip

\noindent Fenster, D.~D.: 2007. Artin in America (1937--1958): A time of transition. In: Reich, K., Kreuzer, A. (Eds.), \emph{Emil Artin (1898--1962): Beitr\"age zu Leben, Werk und Pers\"onlichkeit}. Dr.~Erwin Rauner Verlag, Augsburg, pp.~99--118.\par\medskip

\noindent Frei, G\"unther: The reciprocity law from Euler to Eisenstein. \emph{The intersection of history and mathematics}, 67--90, Sci.~Networks Hist.~Stud., \textbf{15}, Birkh\"auser, Basel, 1994. MR1308080.\par\medskip

\noindent Frei, G\"unther: On the history of the Artin reciprocity law in abelian extensions of algebraic number fields: how Artin was led to his reciprocity law. \emph{The legacy of Niels Henrik Abel}, 267--294, Springer, Berlin, 2004.\par\medskip

\noindent Hasse, Helmut: History of class field theory. 1967 \emph{Algebraic Number Theory} (Proc.~Instructional Conf., Brighton, 1965) pp.~266--279 Thompson, Washington, D.C. MR0218330.\par\medskip

\noindent Herbrand, J.: \emph{Le D\'eveloppment Moderne de la Th\'eorie des Corps Alg\'ebriques: Corps de Classes et Lois de R\'eciprocit\'e}. Gauthier-Villars, Paris, 1936.\par\medskip

\noindent Iyanaga, Sh\=okichi: \emph{Collected papers}. With a survey of Iyanaga's work on number theory by Ichir\=o Satake. Iwanami Shoten, Tokyo, 1994. xx+363~pp. MR1272246.\par\medskip

\noindent Iyanaga, Shokichi: Travaux de Claude Chevalley sur la th\'eorie du corps de classes: introduction. \emph{Jpn. J. Math.} \textbf{1} (2006), no.~1, 25--85. MR2261061.\par\medskip

\noindent Jacobson, Nathan: Abraham Adrian Albert (1905--1972). \emph{Bull. Amer. Math. Soc.} \textbf{80} (1974), 1075--1100. MR0342341.\par\medskip

\noindent Kisilevsky, H.: Olga Taussky-Todd's work in class field theory. \emph{Olga Taussky-Todd: in memoriam}. \emph{Pacific J. Math.} 1997, Special Issue, 219--224.\par\medskip

\noindent Lemmermeyer, Franz: \emph{Reciprocity laws. From Euler to Eisenstein}. Springer Monographs in Mathematics. Springer-Verlag, Berlin, 2000. xx+487~pp.\par\medskip

\noindent Lewis, David W.: Quaternion algebras and the algebraic legacy of Hamilton's quaternions. \emph{Irish Math. Soc. Bull.} No.~57 (2006), 41--64.\par\medskip

\noindent Lorenz, Falko; Roquette, Peter: The theorem of Grunwald--Wang in the setting of valuation theory. \emph{Valuation theory and its applications}, Vol.~II (Saskatoon, SK, 1999), 175--212, Fields Inst.~Commun., \textbf{33}, Amer.~Math.~Soc., Providence, RI, 2003.\par\medskip

\noindent Mac Lane, Saunders: Origins of the cohomology of groups. \emph{Enseign. Math.} (2) \textbf{24} (1978), no.~1--2, 1--29. MR0497280.\par\medskip

\noindent Milne, J.~S.: The Work of John Tate. In Helge Holden and Ragni Piene, editors, \emph{The Abel Prize 2008--2012}, pages 259--340. Springer, Heidelberg, 2014.\par\medskip

\noindent Miyake, Katsuya: The establishment of the Takagi--Artin class field theory. \emph{The intersection of history and mathematics}, 109--128, Sci.~Networks Hist.~Stud., \textbf{15}, Birkh\"auser, Basel, 1994. MR1308082.\par\medskip

\noindent \emph{Class field theory---its centenary and prospect}. Papers from the 7th International Research Institute of the Mathematical Society of Japan (MSJ) held in Tokyo, June 3--12, 1998. Edited by Katsuya Miyake. Advanced Studies in Pure Mathematics, \textbf{30}. Mathematical Society of Japan, Tokyo, 2001. xii+631~pp.\par\medskip

\noindent Miyake, Katsuya: Teiji Takagi, founder of the Japanese School of Modern Mathematics. \emph{Jpn. J. Math.} \textbf{2} (2007), no.~1, 151--164.\par\medskip

\noindent Parshall, Karen Hunger: Joseph H.~M. Wedderburn and the structure theory of algebras. \emph{Arch. Hist. Exact Sci.} \textbf{32} (1985), no.~3--4, 223--349.\par\medskip

\noindent Petri, Birgit; Schappacher, Norbert: From Abel to Kronecker: episodes from 19th century algebra. \emph{The legacy of Niels Henrik Abel}, 227--266, Springer, Berlin, 2004.\par\medskip

\noindent Roquette, Peter: Class field theory in characteristic $p$, its origin and development. \emph{Class field theory---its centenary and prospect} (Tokyo, 1998), 549--631, Adv.~Stud.~Pure Math., \textbf{30}, Math.~Soc.~Japan, Tokyo, 2001.\par\medskip

\noindent Roquette, Peter: \emph{The Brauer--Hasse--Noether theorem in historical perspective}. Schriften der Mathematisch-Naturwissenschaftlichen Klasse der Heidelberger Akademie der Wissenschaften [Publications of the Mathematics and Natural Sciences Section of Heidelberg Academy of Sciences], \textbf{15}. Springer-Verlag, Berlin, 2005. vi+92~pp. ISBN: 3-540-23005-X.\par\medskip

\noindent Schappacher, Norbert: On the history of Hilbert's twelfth problem: a comedy of errors. \emph{Mat\'eriaux pour l'histoire des math\'ematiques au XX\textsuperscript{e} si\`ecle} (Nice, 1996), 243--273, S\'emin.~Congr., \textbf{3}, Soc.~Math.~France, Paris, 1998.\par\medskip

\noindent Snyder, Noah: Artin's $L$-functions: A Historical Approach. Undergrad thesis (B.~Gross). Unpublished.\par\medskip

\noindent Stevenhagen, P.; Lenstra, H.~W.: Chebotar\'ev and his density theorem. \emph{Math. Intelligencer} \textbf{18} (1996), no.~2, 26--37. MR1395088.\par\medskip

\noindent Stevenhagen, Peter: Hilbert's 12th problem, complex multiplication and Shimura reciprocity. \emph{Class field theory---its centenary and prospect} (Tokyo, 1998), 161--176, Adv.~Stud.~Pure Math., \textbf{30}, Math.~Soc.~Japan, Tokyo, 2001.\par\medskip

\noindent Tate, J.: Problem 9: The general reciprocity law. \emph{Mathematical developments arising from Hilbert problems} (Proc.~Sympos.~Pure Math., Northern Illinois Univ., De Kalb, Ill., 1974), pp.~311--322. Proc.~Sympos.~Pure Math., Vol.~XXVIII, Amer.~Math.~Soc., Providence, R.I., 1976. MR0429839.\par\medskip

\noindent Tate, John: Number theory in the 20th century: Part 1. \emph{Bull. Amer. Math. Soc.} (N.S.) \textbf{54} (2017), no.~4, 547--550. MR3683622.
